% =====================================================================
%  RAPPORT DE PROJET — Medicacom AOI Transformation
%
%  Compilation :  latexmk -pdf main.tex
%             ou :  pdflatex main.tex  (deux fois, pour la table des matières)
%
%  Les informations personnelles se règlent ci-dessous, en un seul endroit.
% =====================================================================
\documentclass[12pt,a4paper,openany]{report}

% =====================================================================
%  Préambule — Medicacom AOI Transformation
%  Compilation : pdflatex (deux passes) ou latexmk -pdf main.tex
% =====================================================================

% ---------------------------------------------------------- encodage & langue
\usepackage[T1]{fontenc}
\usepackage[utf8]{inputenc}
\usepackage[french]{babel}
\usepackage{lmodern}
\usepackage{amsmath}
\usepackage{amssymb}
\usepackage{textcomp}
\usepackage{microtype}

% ------------------------------------------------------------------ mise en page
\usepackage[a4paper,left=2.6cm,right=2.4cm,top=2.8cm,bottom=2.6cm]{geometry}
\usepackage{setspace}
\onehalfspacing

% --------------------------------------------------------------------- couleurs
\usepackage[table]{xcolor}
\definecolor{medicaDark}{HTML}{10231C}   % vert sombre de la charte
\definecolor{medicaGreen}{HTML}{157F5C}  % vert accent
\definecolor{medicaDeep}{HTML}{0F5E44}   % vert foncé
\definecolor{medicaTeal}{HTML}{1C6B7D}   % teal (données)
\definecolor{medicaAmber}{HTML}{A66A16}  % ambre (attention)
\definecolor{medicaRed}{HTML}{B23A3A}    % rouge (risque)
\definecolor{medicaSoft}{HTML}{E6F1EB}   % vert très clair
\definecolor{medicaBg}{HTML}{EDF2EF}     % fond
\definecolor{medicaBorder}{HTML}{E1E8E4} % bordure
\definecolor{medicaMuted}{HTML}{5E6F67}  % texte atténué

% ------------------------------------------------------------------ graphiques
\usepackage{graphicx}
\usepackage{float}
\usepackage{tikz}
\usetikzlibrary{
  shapes.geometric, shapes.misc, arrows.meta, positioning, fit, calc,
  backgrounds, decorations.pathreplacing, patterns, matrix, chains
}
\usepackage{pgfplots}
\pgfplotsset{compat=1.18}

% --------------------------------------------------------------------- tableaux
\usepackage{booktabs}
\usepackage{array}
\usepackage{longtable}
\usepackage{tabularx}
\usepackage{multirow}
\usepackage{makecell}
\renewcommand{\arraystretch}{1.28}
\newcolumntype{L}[1]{>{\raggedright\arraybackslash}p{#1}}
\newcolumntype{C}[1]{>{\centering\arraybackslash}p{#1}}
\newcolumntype{R}[1]{>{\raggedleft\arraybackslash}p{#1}}
\newcolumntype{Y}{>{\raggedright\arraybackslash}X}

% ----------------------------------------------------------------------- listes
\usepackage{enumitem}
\setlist{itemsep=2pt, topsep=4pt, parsep=0pt}
\setlist[itemize]{label=\textcolor{medicaGreen}{\textbullet}}

% -------------------------------------------------------------------- légendes
\usepackage{caption}
\usepackage{subcaption}
\captionsetup{
  font=small, labelfont={bf,color=medicaDeep},
  textfont=it, width=0.92\textwidth, skip=7pt
}

% ------------------------------------------------------------------- en-têtes
\usepackage{fancyhdr}
\pagestyle{fancy}
\fancyhf{}
\fancyhead[L]{\small\textcolor{medicaMuted}{\nouppercase{\leftmark}}}
\fancyhead[R]{\small\textcolor{medicaMuted}{Medicacom AOI Transformation}}
\fancyfoot[C]{\small\thepage}
\renewcommand{\headrulewidth}{0.4pt}
\renewcommand{\footrulewidth}{0pt}
\fancypagestyle{plain}{\fancyhf{}\fancyfoot[C]{\small\thepage}%
  \renewcommand{\headrulewidth}{0pt}}

% --------------------------------------------------------------------- titres
\usepackage{titlesec}
\titleformat{\chapter}[display]
  {\normalfont\huge\bfseries\color{medicaDeep}}
  {\filright\Large\textcolor{medicaGreen}{\chaptertitlename\ \thechapter}}
  {8pt}{\Huge}
\titlespacing*{\chapter}{0pt}{-14pt}{26pt}
\titleformat{\section}{\Large\bfseries\color{medicaDeep}}{\thesection}{0.7em}{}
\titleformat{\subsection}{\large\bfseries\color{medicaDark}}{\thesubsection}{0.7em}{}
\titleformat{\subsubsection}{\normalsize\bfseries\color{medicaDark}}{\thesubsubsection}{0.7em}{}

% ------------------------------------------------------------------- encadrés
\usepackage{framed}

% Encadré « Présentation orale suggérée » : sécable sur plusieurs pages.
\newenvironment{oral}{%
  \par\medskip
  \def\FrameCommand{\textcolor{medicaGreen}{\vrule width 3pt}\hspace{12pt}}%
  \MakeFramed{\advance\hsize-\width\FrameRestore}%
  \small\textbf{\textcolor{medicaDeep}{Présentation orale suggérée}}\par\smallskip
  \itshape\ignorespaces
}{\endMakeFramed\par\medskip}

% Encadré neutre pour une remarque importante.
\newenvironment{remarque}[1][Remarque]{%
  \par\medskip
  \def\FrameCommand{\textcolor{medicaTeal}{\vrule width 3pt}\hspace{12pt}}%
  \MakeFramed{\advance\hsize-\width\FrameRestore}%
  \small\textbf{\textcolor{medicaTeal}{#1}}\par\smallskip\ignorespaces
}{\endMakeFramed\par\medskip}

% Encadré d'avertissement, pour les écarts entre l'attendu et le réalisé.
\newenvironment{ecart}[1][Écart avec le cahier des charges]{%
  \par\medskip
  \def\FrameCommand{\textcolor{medicaAmber}{\vrule width 3pt}\hspace{12pt}}%
  \MakeFramed{\advance\hsize-\width\FrameRestore}%
  \small\textbf{\textcolor{medicaAmber}{#1}}\par\smallskip\ignorespaces
}{\endMakeFramed\par\medskip}

% ------------------------------------------------------------------ code source
\usepackage{listings}
\lstdefinestyle{medicacode}{
  basicstyle=\ttfamily\scriptsize,
  keywordstyle=\color{medicaDeep}\bfseries,
  commentstyle=\color{medicaMuted}\itshape,
  stringstyle=\color{medicaTeal},
  numberstyle=\tiny\color{medicaMuted},
  numbers=left, numbersep=8pt,
  backgroundcolor=\color{medicaBg},
  frame=leftline, rulecolor=\color{medicaBorder}, framesep=8pt,
  breaklines=true, breakatwhitespace=false,
  showstringspaces=false, tabsize=2,
  captionpos=b, xleftmargin=14pt,
  extendedchars=true, literate=%
    {é}{{\'e}}1 {è}{{\`e}}1 {ê}{{\^e}}1 {ë}{{\"e}}1
    {à}{{\`a}}1 {â}{{\^a}}1 {î}{{\^\i}}1 {ï}{{\"\i}}1
    {ô}{{\^o}}1 {ö}{{\"o}}1 {û}{{\^u}}1 {ù}{{\`u}}1 {ü}{{\"u}}1
    {ç}{{\c c}}1 {É}{{\'E}}1 {È}{{\`E}}1 {À}{{\`A}}1 {Ç}{{\c C}}1
    {€}{{\texteuro}}1 {×}{{$\times$}}1 {→}{{$\rightarrow$}}1 {−}{{$-$}}1
    {«}{{\guillemotleft}}1 {»}{{\guillemotright}}1 {’}{{'}}1 {…}{{...}}3
}
\lstset{style=medicacode}
\lstdefinelanguage{TypeScript}{
  keywords={const,let,var,function,return,if,else,for,while,import,export,
            from,type,interface,extends,implements,class,new,async,await,
            default,as,of,in,null,undefined,true,false},
  sensitive=true,
  comment=[l]{//}, morecomment=[s]{/*}{*/},
  morestring=[b]", morestring=[b]'
}

% ----------------------------------------------------------------- hyperliens
\usepackage[hidelinks,
  pdftitle={Conception et développement de Medicacom AOI Transformation},
  pdfsubject={Plateforme SaaS de pilotage de l'adoption de l'IA dans l'industrie pharmaceutique},
  pdfkeywords={IA, TOGAF, readiness, portefeuille, pharmaceutique, FastAPI, React}
]{hyperref}
\usepackage[french]{cleveref}

% =====================================================================
%  Commandes utilitaires
% =====================================================================

% Emplacement de capture d'écran. Affiche le vrai fichier s'il existe dans
% images/, sinon un cadre explicite — le document compile dans les deux cas.
\newlength{\captureL}
\newlength{\captureH}
\newcommand{\capture}[3]{%
  \setlength{\captureL}{#3\textwidth}%
  \setlength{\captureH}{0.56\captureL}%
  \begin{figure}[H]
  \centering
  \IfFileExists{images/#1}{%
    \includegraphics[width=\captureL]{images/#1}%
  }{%
    \begin{tikzpicture}
      \draw[dashed, color=medicaBorder, line width=1pt, fill=medicaBg]
        (0pt,0pt) rectangle (\captureL, \captureH);
      \node[align=center, text=medicaMuted, font=\small]
        at (0.5\captureL, 0.5\captureH)
        {\textbf{Capture à insérer}\\[3pt]\ttfamily images/#1};
    \end{tikzpicture}%
  }
  \caption{#2}
  \label{fig:#1}
  \end{figure}%
}

% Marqueurs d'état d'implémentation, employés dans les tableaux de synthèse.
\newcommand{\fait}{\textcolor{medicaGreen}{\textbf{Implémenté}}}
\newcommand{\partiel}{\textcolor{medicaAmber}{\textbf{Partiel}}}
\newcommand{\absent}{\textcolor{medicaRed}{\textbf{Non implémenté}}}
\newcommand{\oui}{\textcolor{medicaGreen}{\textbf{\checkmark}}}
\newcommand{\non}{\textcolor{medicaMuted}{--}}

% Raccourcis de nommage
\newcommand{\appname}{Medicacom AOI Transformation}
\newcommand{\code}[1]{\texttt{\small #1}}

% Styles TikZ partagés par les diagrammes
\tikzset{
  bloc/.style={
    rectangle, rounded corners=4pt, draw=medicaBorder, fill=white,
    text width=2.9cm, align=center, minimum height=1cm, font=\small
  },
  blocvert/.style={bloc, fill=medicaSoft, draw=medicaGreen},
  blocsombre/.style={bloc, fill=medicaDark, draw=medicaDark, text=white},
  blocteal/.style={bloc, fill=white, draw=medicaTeal},
  fleche/.style={-{Stealth[length=5pt]}, draw=medicaMuted, thick},
  flechevert/.style={-{Stealth[length=5pt]}, draw=medicaGreen, thick},
  acteur/.style={
    rectangle, draw=medicaDeep, fill=medicaSoft, rounded corners=3pt,
    text width=2.5cm, align=center, font=\scriptsize\bfseries, minimum height=0.8cm
  },
  cas/.style={
    ellipse, draw=medicaBorder, fill=white, align=center,
    font=\scriptsize, minimum width=2.9cm, minimum height=0.82cm,
    inner sep=2pt
  },
  classe/.style={
    rectangle, draw=medicaDeep, fill=white, align=left,
    font=\scriptsize, inner sep=0pt
  },
  lien/.style={draw=medicaMuted, thin},
}

% En-tête de classe UML : nom + liste d'attributs.
\newcommand{\umlclasse}[2]{%
  \begin{tabular}{@{}l@{}}
    \cellcolor{medicaSoft}\textbf{\strut~#1~}\\
    \hline
    \begin{minipage}{\dimexpr\linewidth}\vspace{2pt}#2\vspace{2pt}\end{minipage}
  \end{tabular}%
}


% ---------------------------------------------------------------------
%  À COMPLÉTER — variables de la page de garde
% ---------------------------------------------------------------------
\newcommand{\varUniversite}{[À COMPLÉTER : nom de l'université]}
\newcommand{\varEcole}{[À COMPLÉTER : nom de l'école ou de la faculté]}
\newcommand{\varDepartement}{[À COMPLÉTER : département]}
\newcommand{\varFormation}{[À COMPLÉTER : intitulé de la formation]}
\newcommand{\varDiplome}{[À COMPLÉTER : diplôme préparé]}

\newcommand{\varTitre}{Conception et développement de Medicacom AOI Transformation}
\newcommand{\varSousTitre}{Plateforme SaaS de pilotage de l'adoption de
  l'Intelligence Artificielle dans l'industrie pharmaceutique}

\newcommand{\varAuteur}{[À COMPLÉTER : prénom et nom]}
\newcommand{\varEncadrantAcademique}{[À COMPLÉTER : encadrant académique]}
\newcommand{\varEncadrantEntreprise}{[À COMPLÉTER : encadrant en entreprise]}
\newcommand{\varEntreprise}{Medicacom}
\newcommand{\varAnnee}{2026 / 2027}

\begin{document}

% =====================================================================
%  Pages liminaires
% =====================================================================
\pagenumbering{roman}
% =====================================================================
%  Page de garde, remerciements, résumé, tables
% =====================================================================

\begin{titlepage}
\begin{tikzpicture}[remember picture, overlay]
  % Bandeau supérieur
  \fill[medicaDark] (current page.north west) rectangle
    ([yshift=-4.2cm]current page.north east);
  % Liseré d'accent
  \fill[medicaGreen] ([yshift=-4.2cm]current page.north west) rectangle
    ([yshift=-4.35cm]current page.north east);
  % Bandeau inférieur
  \fill[medicaDark] ([yshift=1.5cm]current page.south west) rectangle
    (current page.south east);

  % Marque
  \node[anchor=north west, text=white] at ([shift={(2.6cm,-1.5cm)}]current page.north west)
    {\fontsize{26}{28}\selectfont\bfseries Medicacom};
  \node[anchor=north west, text=medicaSoft] at ([shift={(2.6cm,-2.35cm)}]current page.north west)
    {\small\sffamily AOI TRANSFORMATION};

  % Croix pharmaceutique stylisée
  \begin{scope}[shift={([shift={(-3.4cm,-2.1cm)}]current page.north east)}]
    \fill[medicaGreen, rounded corners=6pt] (-0.9,-0.9) rectangle (0.9,0.9);
    \fill[white] (-0.18,-0.6) rectangle (0.18,0.6);
    \fill[white] (-0.6,-0.18) rectangle (0.6,0.18);
  \end{scope}
\end{tikzpicture}

\vspace*{5.2cm}

\begin{center}
  {\large \varUniversite}\\[3pt]
  {\normalsize \varEcole}\\[3pt]
  {\small \textcolor{medicaMuted}{\varDepartement}}

  \vspace{1.6cm}

  {\small\textcolor{medicaMuted}{\textsc{Rapport de projet}}}\\[2pt]
  {\small\textcolor{medicaMuted}{\varDiplome{} --- \varFormation}}

  \vspace{1.3cm}

  \rule{0.72\textwidth}{0.8pt}

  \vspace{0.9cm}

  {\LARGE\bfseries\color{medicaDeep} \varTitre\par}

  \vspace{0.7cm}

  {\large\itshape\color{medicaMuted} \varSousTitre\par}

  \vspace{0.9cm}

  \rule{0.72\textwidth}{0.8pt}

  \vspace{2.1cm}

  \begin{minipage}{0.44\textwidth}
    \raggedright
    \textbf{Réalisé par}\\[4pt]
    \varAuteur
  \end{minipage}
  \hfill
  \begin{minipage}{0.44\textwidth}
    \raggedleft
    \textbf{Encadré par}\\[4pt]
    \varEncadrantAcademique\\
    \varEncadrantEntreprise
  \end{minipage}

  \vspace{1.5cm}

  \textbf{Entreprise d'accueil :} \varEntreprise

  \vfill

  {\color{white}\textbf{Année universitaire \varAnnee}}
  \vspace{0.35cm}
\end{center}
\end{titlepage}

% ---------------------------------------------------------------- remerciements
\chapter*{Remerciements}
\addcontentsline{toc}{chapter}{Remerciements}
\markboth{Remerciements}{}

Je tiens à remercier \varEncadrantAcademique, dont l'encadrement académique a
structuré la conduite de ce travail, ainsi que \varEncadrantEntreprise, pour
l'accompagnement quotidien au sein de \varEntreprise{} et la confiance accordée sur un
sujet aussi transverse que le pilotage de l'adoption de l'intelligence artificielle.

Mes remerciements vont également aux équipes métier qui ont accepté de décrire leurs
processus, leurs irritants et leurs indicateurs. Sans ces échanges, le référentiel de
diagnostic et le questionnaire de qualification n'auraient été qu'un exercice
théorique.

Je remercie enfin les membres du jury pour le temps consacré à l'évaluation de ce
travail.

% ---------------------------------------------------------------------- résumé
\chapter*{Résumé}
\addcontentsline{toc}{chapter}{Résumé}
\markboth{Résumé}{}

Les laboratoires pharmaceutiques accumulent des idées d'intelligence artificielle sans
disposer d'un cadre commun pour décider lesquelles financer. Les besoins remontent de
départements cloisonnés, les données ne sont pas toujours exploitables, la conformité
réglementaire pèse sur chaque décision, et la valeur créée reste rarement mesurée après
la mise en production.

Ce projet répond à ce problème par une plateforme SaaS multi-entreprises,
\appname, qui couvre l'ensemble du parcours : diagnostic de maturité inspiré de
la démarche TOGAF, qualification de la solution la plus adaptée, évaluation du
readiness, priorisation du portefeuille, génération de la fiche projet, suivi du cycle
de vie et mesure du retour sur investissement.

La plateforme repose sur une API REST écrite en Python avec FastAPI et SQLAlchemy~2,
une base PostgreSQL et une interface React en TypeScript. Le contrôle d'accès associe
41~permissions granulaires, 11~rôles fonctionnels et cinq niveaux de périmètre, avec
une isolation stricte entre organisations clientes. Chaque décision produite par le
système est explicable : le score de priorité renvoie la contribution de chacun de ses
six critères, et la recommandation de solution expose les règles satisfaites, les
alternatives écartées et les garde-fous à mettre en place.

Un point mérite d'être énoncé sans détour : la plateforme ne contient aujourd'hui
\emph{aucun modèle d'apprentissage automatique}. Son intelligence est celle d'un moteur
de règles pondérées, déterministe et auditable. Ce choix est assumé et documenté au
chapitre~\ref{chap:realisation}; l'intégration d'un modèle de langage figure parmi les
perspectives, non parmi les fonctionnalités livrées.

\vspace{0.8em}
\noindent\textbf{Mots-clés} — adoption de l'IA, industrie pharmaceutique, TOGAF,
readiness, portefeuille de projets, moteur de règles, gouvernance, multi-tenant,
FastAPI, React, PostgreSQL.

% -------------------------------------------------------------------- abstract
\chapter*{Abstract}
\addcontentsline{toc}{chapter}{Abstract}
\markboth{Abstract}{}

Pharmaceutical companies accumulate artificial intelligence ideas without a shared
framework for deciding which ones deserve funding. Needs come from siloed departments,
data is not always usable, regulatory compliance weighs on every decision, and the
value actually delivered is rarely measured once a system reaches production.

This project addresses that problem with \appname, a multi-tenant SaaS platform
covering the full journey: a TOGAF-inspired maturity assessment, qualification of the
most suitable solution type, readiness scoring, portfolio prioritisation, project
charter generation, lifecycle tracking and return-on-investment measurement.

The platform is built on a REST API written in Python with FastAPI and SQLAlchemy~2, a
PostgreSQL database and a React frontend in TypeScript. Access control combines
41~granular permissions, 11~functional roles and five scope levels, with strict
isolation between client organisations. Every decision the system produces is
explainable: the priority score returns the contribution of each of its six criteria,
and the solution recommendation exposes the rules that fired, the alternatives that
were ruled out and the guardrails to put in place.

One point deserves to be stated plainly: the platform currently contains
\emph{no machine learning model}. Its intelligence is that of a weighted, deterministic
and auditable rules engine. This is a deliberate choice, documented in
Chapter~\ref{chap:realisation}; integrating a large language model belongs to the
outlook, not to the delivered features.

\vspace{0.8em}
\noindent\textbf{Keywords} — AI adoption, pharmaceutical industry, TOGAF, readiness,
project portfolio, rules engine, governance, multi-tenancy, FastAPI, React, PostgreSQL.

% ----------------------------------------------------------------------- tables
\cleardoublepage
\renewcommand{\contentsname}{Table des matières}
\tableofcontents

\cleardoublepage
\listoffigures
\addcontentsline{toc}{chapter}{Table des figures}

\cleardoublepage
\listoftables
\addcontentsline{toc}{chapter}{Liste des tableaux}


% =====================================================================
%  Corps du rapport
% =====================================================================
\cleardoublepage
\pagenumbering{arabic}

% =====================================================================
%  Introduction générale
% =====================================================================
\chapter*{Introduction générale}
\addcontentsline{toc}{chapter}{Introduction générale}
\markboth{Introduction générale}{}
\label{chap:introduction}

\section*{Un secteur sous double contrainte}

L'industrie pharmaceutique se transforme sous deux pressions simultanées. D'un côté,
la pression économique : les coûts de développement d'un médicament augmentent, les
marges se resserrent, et l'optimisation des opérations devient un levier de
compétitivité. De l'autre, la pression réglementaire : chaque processus touchant au
produit, au lot ou au dossier patient relève d'un cadre GxP qui impose traçabilité,
validation informatique et justification des décisions.

L'intelligence artificielle est présentée comme une réponse à la première pression.
Elle se heurte pourtant frontalement à la seconde. Un modèle dont on ne peut pas
expliquer la sortie est difficilement défendable devant un inspecteur. Une
automatisation qui décide seule sur un périmètre GxP est encore plus difficile à
justifier.

\section*{Le problème observé}

Dans les organisations pharmaceutiques, l'adoption de l'IA échoue rarement par manque
d'idées. Elle échoue plus souvent en amont du code, pour des raisons qui n'ont rien de
technique :

\begin{itemize}
  \item \textbf{les besoins remontent de départements cloisonnés.} Le Commercial, la
    Production, la Qualité et les Affaires réglementaires expriment leurs irritants
    chacun de leur côté, souvent dans des formats différents, parfois sur le même
    processus sans le savoir ;
  \item \textbf{la qualification technique arrive trop tard.} On décide de « faire de
    l'IA » avant de vérifier si les données existent, si elles sont exploitables et si
    une vérité terrain est disponible. Le projet s'arrête alors au premier atelier
    technique ;
  \item \textbf{les arbitrages manquent de critères partagés.} Sans grille commune,
    la priorisation se joue à la conviction du sponsor le plus influent plutôt qu'à
    l'impact attendu ;
  \item \textbf{la valeur n'est pas mesurée.} Le business case sert à obtenir le budget,
    puis disparaît. Personne ne compare la valeur promise à la valeur réalisée ;
  \item \textbf{la traçabilité fait défaut.} Six mois plus tard, plus personne ne sait
    pourquoi telle solution a été retenue plutôt qu'une autre.
\end{itemize}

Ces cinq difficultés partagent une même cause : l'absence d'un référentiel unique où
le diagnostic, la qualification, l'arbitrage et le suivi vivent ensemble, avec leurs
justifications.

\section*{La réponse apportée}

\appname{} est une plateforme SaaS multi-entreprises qui couvre le parcours complet,
de l'état des lieux à la valeur captée. Elle s'organise autour de sept temps
(figure~\ref{fig:parcours}).

\begin{figure}[H]
\centering
\begin{tikzpicture}[node distance=0.42cm and 0.5cm]
  \node[blocvert] (diag) {\textbf{Diagnostic}\\\scriptsize AS-IS, maturité};
  \node[blocvert, right=of diag] (qual) {\textbf{Qualification}\\\scriptsize type de solution};
  \node[blocvert, right=of qual] (read) {\textbf{Readiness}\\\scriptsize score sur 100};
  \node[blocvert, below=0.9cm of read] (prio) {\textbf{Priorisation}\\\scriptsize score, rang, vague};
  \node[blocvert, left=of prio] (fiche) {\textbf{Fiche projet}\\\scriptsize versions, décision};
  \node[blocvert, left=of fiche] (cycle) {\textbf{Cycle de vie}\\\scriptsize phases, gates};
  \node[blocsombre, below=0.9cm of fiche] (roi) {\textbf{Valeur \& ROI}\\\scriptsize TCO, gains, payback};

  \draw[flechevert] (diag) -- (qual);
  \draw[flechevert] (qual) -- (read);
  \draw[flechevert] (read) -- (prio);
  \draw[flechevert] (prio) -- (fiche);
  \draw[flechevert] (fiche) -- (cycle);
  \draw[flechevert] (cycle.south) |- (roi.west);
  \draw[fleche, dashed] (roi.east) -| node[pos=0.25, below, font=\scriptsize, text=medicaMuted]
    {rétroaction : la valeur mesurée corrige les hypothèses} (diag.south);
\end{tikzpicture}
\caption{Le parcours couvert par la plateforme, du diagnostic à la mesure de la valeur}
\label{fig:parcours}
\end{figure}

Chacun de ces temps produit une décision, et chaque décision reste rattachée à ce qui
l'a motivée. Le readiness renvoie la note de chaque dimension avec sa preuve. La
recommandation de solution expose les règles satisfaites et les alternatives écartées.
Le score de priorité détaille la contribution de ses six critères. Le business case
publie ses coûts et ses gains ligne à ligne.

\section*{Ce que ce rapport établit}

Ce document décrit ce qui a été \emph{réellement} développé. Lorsque l'implémentation
s'écarte du cahier des charges initial, l'écart est signalé explicitement dans un
encadré dédié, avec ce qui était demandé, ce qui a été réalisé et ce qui reste à faire.

Un point mérite d'être posé dès l'introduction, car il conditionne la lecture de tout
le rapport. Le nom du produit contient « IA », et la plateforme recommande des types
de solutions d'intelligence artificielle. Elle ne \emph{contient} pourtant aucun modèle
d'apprentissage automatique : ni modèle entraîné localement, ni appel à un fournisseur
externe. Sa capacité de recommandation repose sur un moteur de règles pondérées,
déterministe et explicable. Ce choix est argumenté à la
section~\ref{sec:pas-dia}~; le confondre avec de l'apprentissage automatique serait une
erreur d'analyse.

\section*{Organisation du document}

Le rapport suit la progression du projet.

\begin{itemize}
  \item Le \textbf{chapitre~\ref{chap:analyse}} établit le contexte, la problématique,
    les objectifs, le cahier des charges et l'état réel de chacun des douze modules
    fonctionnels.
  \item Le \textbf{chapitre~\ref{chap:conception}} présente les acteurs, le modèle de
    droits, les cas d'utilisation, le modèle de données et l'architecture technique.
  \item Le \textbf{chapitre~\ref{chap:realisation}} détaille les technologies employées,
    les algorithmes de readiness, de priorisation et de valorisation, le moteur de
    qualification et le dispositif de sécurité.
  \item Le \textbf{chapitre~\ref{chap:interfaces}} parcourt les quatorze écrans de
    l'application. Chaque écran est accompagné d'une proposition de présentation orale
    directement réutilisable en soutenance.
  \item Le \textbf{chapitre~\ref{chap:tests}} rend compte de la stratégie de validation
    et des résultats obtenus.
  \item Le \textbf{chapitre~\ref{chap:conclusion}} dresse le bilan, énonce les limites
    de la version actuelle et ouvre sur les évolutions envisageables.
\end{itemize}

L'annexe~\ref{ann:soutenance} propose un guide de démonstration minuté et une série de
questions susceptibles d'être posées par le jury, avec des réponses appuyées sur le
code réel.

% =====================================================================
%  Chapitre 1 — Étude et analyse des besoins
% =====================================================================
\chapter{Étude et analyse des besoins}
\label{chap:analyse}

\section{Contexte du projet}

Le projet se déroule chez \varEntreprise, laboratoire pharmaceutique dont le jeu de
données de démonstration reflète l'organisation : quatre implantations — un site de
production à Tunis, un centre de R\&D à Lyon, une plateforme de distribution à
Casablanca et un siège à Paris — et huit départements, du Commercial aux Affaires
réglementaires.

Cette dispersion géographique et fonctionnelle est précisément ce qui rend
l'arbitrage difficile. Un besoin de prévision des ventes exprimé par le Commercial de
Tunis et un besoin de prévision de la demande formulé par la Supply Chain de
Casablanca reposent sur le même référentiel produit, mais rien ne le signale tant que
les deux dossiers vivent dans des tableurs distincts.

\section{Présentation de \appname}

\appname{} est une plateforme de pilotage de portefeuille, destinée à la direction, aux
responsables de département et aux fonctions transverses d'un laboratoire. Elle n'est
pas un outil de développement de modèles : elle intervient \emph{avant} le projet
technique, pour décider s'il doit exister, et \emph{après} sa mise en production, pour
vérifier qu'il tient ses promesses.

La plateforme est multi-entreprises. Chaque organisation cliente dispose de ses sites,
départements, utilisateurs, rôles et données métier, sans aucune visibilité sur les
autres.

\section{Problématique}

La question à laquelle le projet répond peut s'énoncer ainsi :

\begin{remarque}[Problématique]
Comment doter une organisation pharmaceutique d'un cadre unique et traçable qui
transforme des idées d'intelligence artificielle dispersées en un portefeuille de
projets arbitrés, gouvernés et dont la valeur est mesurée, tout en respectant les
exigences d'explicabilité et de conformité propres au secteur ?
\end{remarque}

Trois sous-questions en découlent :

\begin{enumerate}
  \item Comment qualifier objectivement le type de solution adapté à un besoin, sans
    céder à l'effet de mode qui pousse à répondre « IA » à toute question ?
  \item Comment comparer des opportunités issues de départements différents sur une
    grille commune, tout en laissant les pondérations ajustables selon la stratégie du
    moment ?
  \item Comment garantir qu'une décision prise il y a six mois reste justifiable devant
    un comité ou un inspecteur ?
\end{enumerate}

\section{Objectifs généraux}

\begin{enumerate}
  \item Centraliser le diagnostic de maturité, l'expression des besoins et le
    portefeuille de projets dans un référentiel unique.
  \item Produire des décisions explicables : toute note, tout score et toute
    recommandation doivent pouvoir être décomposés.
  \item Garantir l'étanchéité entre organisations clientes et le cloisonnement des
    droits à l'intérieur de chacune.
  \item Mesurer la valeur réellement captée et la confronter à la valeur promise.
\end{enumerate}

\section{Objectifs spécifiques}

\begin{itemize}
  \item Évaluer la maturité sur deux axes, technique et métier, avec pondération et
    niveau de confiance par dimension.
  \item Identifier les dimensions bloquantes, c'est-à-dire celles qui interdisent un
    passage en production tant qu'elles ne sont pas remédiées.
  \item Qualifier chaque opportunité parmi huit types de solution, avec justification
    écrite, alternatives écartées et garde-fous.
  \item Calculer un score de priorité sur six critères pondérables, diminué d'une
    pénalité de risque.
  \item Détecter les doublons et les synergies entre opportunités de départements
    différents.
  \item Générer une fiche projet versionnée et soumise à décision de comité.
  \item Suivre le cycle de vie par phases avec gates de passage.
  \item Calculer TCO, gains annuels, ROI et délai de retour, au projet et au
    portefeuille.
  \item Journaliser toute modification significative de façon immuable.
\end{itemize}

\section{Besoins fonctionnels}

Les besoins fonctionnels se regroupent en douze modules, notés M1 à M12. Le
tableau~\ref{tab:modules} en donne l'état réel dans l'application livrée. Cet état a
été établi en inspectant le code : routers FastAPI, services métier, écrans React et
suite de tests.

\begin{longtable}{@{}L{1.1cm} L{6.6cm} L{2.6cm} L{4.2cm}@{}}
\caption{État d'implémentation des douze modules fonctionnels}
\label{tab:modules}\\
\toprule
\textbf{Mod.} & \textbf{Description} & \textbf{État} & \textbf{Constat dans le code} \\
\midrule
\endfirsthead
\multicolumn{4}{c}{\small\itshape suite du tableau~\ref{tab:modules}}\\
\toprule
\textbf{Mod.} & \textbf{Description} & \textbf{État} & \textbf{Constat dans le code} \\
\midrule
\endhead
\bottomrule
\endfoot

M1 & Organisation, sites, départements, utilisateurs, rôles et permissions
   & \fait
   & 18 endpoints sous \code{/org}, 41 permissions, 11 rôles, écran Administration à
     trois onglets \\

M2 & Diagnostic AS-IS inspiré de TOGAF : campagnes et réponses pondérées
   & \fait
   & 8 endpoints \code{/diagnostic}, questionnaire éditable dans l'interface \\

M3 & Expression des besoins : questionnaire commun et volet départemental
   & \fait
   & \code{services/questionnaire.py}, 5 sections communes et 7 volets métier \\

M4 & Qualification du type de solution
   & \fait
   & \code{services/recommendation.py}, 8 types, 39 règles pondérées \\

M5 & Readiness : score sur 100, dimensions bloquantes
   & \fait
   & \code{services/readiness.py}, distinction bloquante / vigilance / à instruire \\

M6 & Croisements interservices : processus, applications, données, dépendances,
     doublons
   & \fait
   & 14 endpoints \code{/crossings}, écran Architecture \& données \\

M7 & Priorisation du portefeuille, pondérations paramétrables
   & \fait
   & \code{services/scoring.py}, panneau de repriorisation avec curseurs \\

M8 & Fiche projet : génération, versions, validation, export
   & \partiel
   & Génération, versionnement et décisions opérationnels ; export en JSON et HTML
     imprimable, pas de génération \code{.docx} \\

M9 & Suivi du cycle de vie : phases, jalons, tâches, risques, gates
   & \partiel
   & API complète (16 endpoints \code{/projects}) mais aucun écran de gestion de
     projet dédié \\

M10 & Pilotage de la valeur : TCO, gains, ROI, payback
   & \fait
   & \code{services/value.py}, 9 endpoints \code{/value}, écran Valeur \& ROI \\

M11 & Tableaux de bord et agrégations
   & \fait
   & 4 endpoints \code{/dashboards}, vues Direction et Département \\

M12 & Administration et journal d'audit immuable
   & \partiel
   & Journal complet et écran dédié ; notifications persistées mais sans écran de
     consultation \\
\end{longtable}

\begin{ecart}
Trois modules ne sont que partiellement couverts, et il serait malhonnête de les
présenter autrement.

\textbf{M8 — export.} Le cahier des charges prévoyait un export bureautique. Le bouton
de l'interface porte encore le libellé « Export Word / PDF » alors que l'API produit du
JSON et une page HTML mise en page pour l'impression. L'utilisateur obtient bien un PDF
en passant par la fonction d'impression de son navigateur, mais aucun fichier
\code{.docx} n'est généré. Le libellé est donc en avance sur la fonction.

\textbf{M9 — interface de suivi.} Les endpoints existent et sont protégés : phases,
jalons, tâches, risques, mitigations, incidents et passage de gates. Aucun écran dédié
ne les expose. Le PMO retrouve ses projets, leur avancement et leur prochaine gate dans
le tableau de bord de département, ce qui couvre le pilotage mais pas la gestion
détaillée des tâches.

\textbf{M12 — notifications.} La table et les endpoints existent, l'icône de cloche de
l'en-tête affiche une pastille de non-lu, mais aucun panneau ne permet de lire les
notifications.
\end{ecart}

\section{Besoins non fonctionnels}

\begin{longtable}{@{}L{3.4cm} L{10.8cm}@{}}
\caption{Besoins non fonctionnels et traitement retenu}
\label{tab:nfr}\\
\toprule
\textbf{Exigence} & \textbf{Traitement dans l'application} \\
\midrule
\endfirsthead
\toprule
\textbf{Exigence} & \textbf{Traitement dans l'application} \\
\midrule
\endhead
\bottomrule
\endfoot

Étanchéité des tenants
& Toute table métier porte \code{organization\_id}. Le jeton JWT transporte
  l'identifiant du tenant et la dépendance \code{get\_tenant\_id} le lit depuis le
  jeton, jamais depuis la requête. Neuf tests couvrent ce point. \\

Explicabilité
& Trois endpoints renvoient le raisonnement et non le seul résultat :
  \code{/opportunities/\{id\}/score}, \code{/opportunities/\{id\}/qualify} et
  \code{/diagnostic/readiness}. \\

Traçabilité
& Journal \code{audit\_log} en insertion seule : l'API n'expose ni mise à jour ni
  suppression sur cette table. Chaque entrée conserve auteur, horodatage, entité,
  ancienne et nouvelle valeur, et adresse IP. \\

Sécurité des accès
& Authentification par jeton JWT signé HS256, mots de passe hachés avec bcrypt,
  contrôle de permission sur chaque endpoint sensible. \\

Robustesse du schéma
& Migrations Alembic versionnées, écrites pour s'appliquer sur une base déjà peuplée
  sans perte de données. \\

Maintenabilité
& Logique métier isolée dans sept services sans dépendance à FastAPI, donc testables
  indépendamment du transport HTTP. \\

Ergonomie
& Interface en français, navigation et actions dérivées des permissions réelles,
  responsive jusqu'à 390~pixels de large. \\
\end{longtable}

\section{Contraintes}

\begin{itemize}
  \item \textbf{Contrainte réglementaire.} Le périmètre GxP impose que toute
    recommandation reste explicable et qu'une validation humaine subsiste. Le moteur de
    qualification ajoute automatiquement trois garde-fous lorsque l'opportunité est
    déclarée GxP : validation informatique, traçabilité ALCOA+ et revue humaine
    obligatoire.
  \item \textbf{Contrainte de confidentialité.} Deux laboratoires concurrents peuvent
    être hébergés sur la même instance. L'isolation ne peut pas dépendre du bon
    comportement du frontend.
  \item \textbf{Contrainte de langue.} L'interface et les données de démonstration sont
    en français, ce qui a imposé une attention particulière aux jeux de caractères,
    notamment pour l'export CSV destiné à Excel.
  \item \textbf{Contrainte d'environnement.} Le développement s'est fait sous Windows
    avec PowerShell, ce qui a influencé la documentation d'installation.
\end{itemize}

\section{Méthodologie adoptée}

Le développement a suivi une progression par couches, du socle vers l'interface, avec
validation à chaque étape.

\begin{enumerate}
  \item \textbf{Modélisation des données.} Les 34~tables ont été définies avant tout
    développement d'API, puis figées par une migration Alembic initiale.
  \item \textbf{Services métier.} Readiness, qualification, priorisation et
    valorisation ont été écrits comme des fonctions pures, testables sans base de
    données ni serveur.
  \item \textbf{Exposition HTTP.} Les routers FastAPI se contentent d'orchestrer :
    vérifier les droits, charger les entités, appeler le service, journaliser.
  \item \textbf{Jeu de démonstration.} Un seed réaliste a été construit tôt, pour que
    chaque écran soit développé face à des données crédibles plutôt qu'à des tableaux
    vides.
  \item \textbf{Interface.} Les écrans ont été développés un par un, chacun vérifié par
    capture dans un navigateur réel avant de passer au suivant.
  \item \textbf{Durcissement.} Le contrôle d'accès a été repris en profondeur dans un
    second temps, une fois les fonctionnalités stabilisées, et couvert par une suite de
    tests dédiée.
\end{enumerate}

Cette dernière étape mérite d'être soulignée : la première version du contrôle d'accès
ne comptait que dix permissions, toutes en écriture. Les quarante-six endpoints de
lecture n'étaient protégés que par l'authentification, ce qui laissait tout utilisateur
connecté lire l'intégralité de son organisation. La reprise a porté le catalogue à
41~permissions et introduit la notion de périmètre.

% =====================================================================
%  Chapitre 2 — Conception
% =====================================================================
\chapter{Conception du système}
\label{chap:conception}

\section{Acteurs et rôles}

Le modèle de droits distingue onze rôles fonctionnels. Chacun combine un jeu de
permissions et un niveau de périmètre, c'est-à-dire l'étendue des données sur
lesquelles il agit. Les descriptions qui suivent reprennent le catalogue défini dans
\code{app/core/rbac.py}.

\subsection{Administrateur Medicacom}

\textbf{Mission.} Administrer la plateforme et les organisations clientes : référentiels,
modèles de questionnaires, rôles disponibles et support des tenants.

\textbf{Périmètre.} Plateforme — le seul rôle qui franchit la frontière des
organisations.

\textbf{Droits.} Les 41~permissions du catalogue, dont \code{platform.admin}, réservée à
ce rôle. Un administrateur d'entreprise ne peut ni se l'attribuer ni la placer dans un
rôle qu'il crée.

\subsection{Administrateur entreprise}

\textbf{Mission.} Administrer son organisation : sites, départements, utilisateurs,
rôles et paramètres.

\textbf{Périmètre.} Organisation. Il n'accède jamais à une autre.

\textbf{Droits.} 23~permissions : toutes les lectures, plus la gestion de
l'organisation, des sites, des départements, des comptes et des rôles, ainsi que
l'accès au journal d'audit.

\textbf{Limitations.} Il n'arbitre pas le portefeuille et ne renseigne pas les
diagnostics. Son écran d'accueil est la console d'administration, pas un tableau de
bord.

\subsection{Direction Générale}

\textbf{Mission.} Vision globale du portefeuille, arbitrage et décision de comité.

\textbf{Périmètre.} Organisation.

\textbf{Droits.} 22~permissions : toutes les lectures, le tableau de bord Direction, la
validation d'opportunité, la repriorisation du portefeuille, le passage de gate,
la décision de comité sur la fiche projet et la validation des bénéfices financiers.

\textbf{Limitations.} La Direction n'administre pas les comptes et ne renseigne pas le
questionnaire de diagnostic. Elle le consulte.

\subsection{Responsable de département}

\textbf{Mission.} Piloter les besoins et les projets de son ou ses départements.

\textbf{Périmètre.} Département — le seul rôle dont la portée est restreinte par
rattachement.

\textbf{Droits.} 18~permissions : création et modification d'opportunités, réponse et
validation du diagnostic, rédaction de fiche projet, tableau de bord de département.

\textbf{Limitations.} Il consulte les dépendances transverses mais ne modifie rien hors
de ses départements. Une tentative d'écriture ailleurs répond 403.

\subsection{Contributeur métier}

\textbf{Mission.} Répondre aux questionnaires assignés, documenter les processus et
proposer des opportunités.

\textbf{Périmètre.} Contributions personnelles.

\textbf{Droits.} 11~permissions, le jeu le plus restreint : lecture de l'organisation et
de son département, réponse au diagnostic, création et lecture d'opportunités, lecture
des recommandations et des fiches.

\textbf{Limitations.} Aucun accès au portefeuille consolidé ni aux données financières.
Aucun arbitrage, aucune validation.

\subsection{DSI / Architecte SI}

\textbf{Mission.} Évaluer la faisabilité technique, tenir l'architecture AS-IS et TO-BE
ainsi que les dépendances applicatives.

\textbf{Droits.} 16~permissions, dont \code{architecture.manage} qui lui est propre avec
l'administrateur, et \code{project.validate\_gate} pour les gates techniques.

\subsection{Data Office / Data Engineer}

\textbf{Mission.} Tenir le référentiel des sources de données, leur qualité, leur
ownership et les dépendances data interservices.

\textbf{Droits.} 14~permissions, dont \code{data.manage}. C'est le seul rôle, avec
l'administrateur de plateforme, qui écrit dans le référentiel des sources.

\textbf{Limitations.} Il ne gère pas le référentiel applicatif, qui relève du DSI.

\subsection{AI Lead / Data Scientist}

\textbf{Mission.} Analyser les opportunités qualifiées, amender la recommandation avec
justification, définir les critères de performance et les garde-fous.

\textbf{Droits.} 16~permissions, dont \code{recommendation.override}, qui lui est
exclusive. C'est le seul rôle autorisé à relancer la qualification d'une opportunité.

\textbf{Limitations.} Il n'administre pas les comptes.

\subsection{Finance / Contrôle de gestion}

\textbf{Mission.} Instruire CAPEX, OPEX et gains, valider les hypothèses du business
case, comparer valeur attendue et valeur réalisée.

\textbf{Droits.} 15~permissions, dont \code{finance.manage} et \code{finance.validate}.

\textbf{Limitations.} Finance ne touche ni à la recommandation IA ni à l'architecture
technique. Un appel à l'endpoint de qualification depuis un compte Finance répond 403.

\subsection{PMO / Chef de projet}

\textbf{Mission.} Gérer le cycle de vie des projets : phases, jalons, tâches, risques,
incidents, budget de suivi et reporting.

\textbf{Droits.} 15~permissions, dont \code{project.manage}.

\textbf{Limitations.} Il ne valide pas les bénéfices financiers.

\subsection{Auditeur / Lecture seule}

\textbf{Mission.} Consulter preuves, scores, versions, décisions et journal d'audit.

\textbf{Droits.} 16~permissions, toutes en lecture, plus \code{audit.read}.

\textbf{Limitations.} Aucune création, modification, suppression ni validation, quel que
soit le module. Huit tests paramétrés vérifient ce point endpoint par endpoint.

\subsection{Matrice de synthèse}

\begin{table}[H]
\centering
\caption{Capacités par rôle, dérivées du catalogue de permissions}
\label{tab:matrice-roles}
\small
\begin{tabular}{@{}L{4.4cm} C{1.6cm} C{1.5cm} C{1.8cm} C{1.7cm} C{1.9cm}@{}}
\toprule
\textbf{Rôle} & \textbf{Consul\-tation} & \textbf{Créa\-tion} &
\textbf{Modifi\-cation} & \textbf{Valida\-tion} & \textbf{Adminis\-tration} \\
\midrule
Administrateur Medicacom   & \oui & \oui & \oui & \oui & \oui \\
Administrateur entreprise  & \oui & \oui & \oui & \non & \oui \\
Direction Générale         & \oui & \non & \oui & \oui & \non \\
Responsable de département & \oui & \oui & \oui & \oui & \non \\
Contributeur métier        & \oui & \oui & \oui & \non & \non \\
DSI / Architecte SI        & \oui & \oui & \oui & \oui & \non \\
Data Office / Data Engineer& \oui & \oui & \oui & \non & \non \\
AI Lead / Data Scientist   & \oui & \oui & \oui & \oui & \non \\
Finance / Contrôle gestion & \oui & \oui & \oui & \oui & \non \\
PMO / Chef de projet       & \oui & \oui & \oui & \non & \non \\
Auditeur / Lecture seule   & \oui & \non & \non & \non & \non \\
\bottomrule
\end{tabular}
\end{table}

La colonne \emph{Validation} recouvre des objets différents selon le rôle, ce que le
tableau ne peut pas montrer. La Direction valide les opportunités, les gates, les
fiches et les bénéfices. Le responsable de département ne valide que les campagnes de
diagnostic. Le DSI et l'AI Lead ne valident que les gates de projet. Finance ne valide
que les bénéfices financiers. La granularité réelle est celle de la permission, pas
celle de la colonne.

\section{Cas d'utilisation}

\subsection{Vue générale}

\begin{figure}[H]
\centering
\begin{tikzpicture}[scale=0.92, transform shape]
  % --- acteurs à gauche
  \node[acteur] (dir)   at (0,4.6)  {Direction Générale};
  \node[acteur] (resp)  at (0,3.0)  {Responsable de département};
  \node[acteur] (contrib) at (0,1.4) {Contributeur métier};
  \node[acteur] (adm)   at (0,-0.2) {Administrateur entreprise};

  % --- acteurs à droite
  \node[acteur] (dsi)   at (12,4.6) {DSI / Architecte};
  \node[acteur] (data)  at (12,3.0) {Data Office};
  \node[acteur] (ai)    at (12,1.4) {AI Lead};
  \node[acteur] (fin)   at (12,-0.2){Finance};
  \node[acteur] (pmo)   at (12,-1.8){PMO};
  \node[acteur] (audit) at (0,-1.8) {Auditeur};

  % --- système
  \begin{scope}[on background layer]
    \draw[draw=medicaBorder, fill=medicaBg, rounded corners=8pt]
      (2.7,-2.9) rectangle (9.3,5.6);
  \end{scope}
  \node[font=\small\bfseries, text=medicaDeep] at (6,5.25) {Medicacom AOI Transformation};

  \node[cas] (connect) at (6,4.45) {Se connecter};
  \node[cas] (diag)    at (4.4,3.5) {Réaliser le diagnostic};
  \node[cas] (repond)  at (7.6,3.5) {Répondre au questionnaire};
  \node[cas] (opp)     at (4.4,2.5) {Créer une opportunité};
  \node[cas] (qualif)  at (7.6,2.5) {Qualifier l'opportunité};
  \node[cas] (reco)    at (4.4,1.5) {Consulter la recommandation};
  \node[cas] (over)    at (7.6,1.5) {Amender la recommandation};
  \node[cas] (prio)    at (4.4,0.5) {Prioriser le portefeuille};
  \node[cas] (fiche)   at (7.6,0.5) {Générer la fiche projet};
  \node[cas] (valid)   at (4.4,-0.5) {Valider en comité};
  \node[cas] (suivi)   at (7.6,-0.5) {Suivre le projet};
  \node[cas] (roi)     at (4.4,-1.5) {Consulter le ROI};
  \node[cas] (admin)   at (7.6,-1.5) {Administrer l'organisation};
  \node[cas] (journal) at (6,-2.4) {Consulter le journal d'audit};

  % --- associations principales
  \foreach \a in {dir,resp,contrib,adm,audit}
    \draw[lien] (\a.east) -- (connect.west);
  \foreach \a in {dsi,data,ai,fin,pmo}
    \draw[lien] (\a.west) -- (connect.east);

  \draw[lien] (dir.east)   -- (prio.west);
  \draw[lien] (dir.east)   -- (valid.west);
  \draw[lien] (dir.east)   -- (roi.west);
  \draw[lien] (resp.east)  -- (opp.west);
  \draw[lien] (resp.east)  -- (diag.west);
  \draw[lien] (contrib.east) -- (repond.west);
  \draw[lien] (contrib.east) -- (opp.west);
  \draw[lien] (adm.east)   -- (admin.west);
  \draw[lien] (audit.east) -- (journal.west);

  \draw[lien] (dsi.west)   -- (qualif.east);
  \draw[lien] (data.west)  -- (repond.east);
  \draw[lien] (ai.west)    -- (over.east);
  \draw[lien] (ai.west)    -- (reco.east);
  \draw[lien] (fin.west)   -- (roi.east);
  \draw[lien] (pmo.west)   -- (suivi.east);
  \draw[lien] (pmo.west)   -- (fiche.east);

  % --- dépendance interne
  \draw[-{Stealth[length=4pt]}, dashed, draw=medicaGreen] (opp) -- node[midway, above, font=\tiny, text=medicaGreen] {\guillemotleft include\guillemotright} (qualif);
\end{tikzpicture}
\caption{Diagramme de cas d'utilisation général}
\label{fig:uc-general}
\end{figure}

La relation \guillemotleft~include~\guillemotright{} entre la création d'une opportunité
et sa qualification n'est pas décorative : dans le code, l'endpoint de création appelle
systématiquement le moteur de qualification avant de répondre. Il n'existe pas
d'opportunité non qualifiée.

\subsection{Sous-système Diagnostic et Readiness}

\begin{figure}[H]
\centering
\begin{tikzpicture}[scale=0.95, transform shape]
  \node[acteur] (resp) at (0,2.2) {Responsable de département};
  \node[acteur] (dsi)  at (0,0.8) {DSI / Data Office};
  \node[acteur] (dir)  at (0,-0.6) {Direction};

  \begin{scope}[on background layer]
    \draw[draw=medicaBorder, fill=medicaBg, rounded corners=8pt] (2.8,-1.6) rectangle (10.4,3.1);
  \end{scope}
  \node[font=\small\bfseries, text=medicaDeep] at (6.6,2.75) {Module M2 / M5};

  \node[cas] (camp)  at (4.6,2.0) {Ouvrir une campagne};
  \node[cas] (note)  at (4.6,1.0) {Noter une dimension};
  \node[cas] (preuve)at (4.6,0.0) {Joindre une preuve};
  \node[cas] (calc)  at (8.4,1.5) {Calculer le readiness};
  \node[cas] (bloq)  at (8.4,0.5) {Identifier les dimensions bloquantes};
  \node[cas] (exp)   at (8.4,-0.6) {Exporter le diagnostic};
  \node[cas] (valcamp) at (4.6,-1.0) {Valider la campagne};

  \draw[lien] (resp.east) -- (camp.west);
  \draw[lien] (resp.east) -- (valcamp.west);
  \draw[lien] (dsi.east)  -- (note.west);
  \draw[lien] (dsi.east)  -- (preuve.west);
  \draw[lien] (dir.east)  -- (exp.west);
  \draw[lien] (dir.east)  -- (calc.west);

  \draw[-{Stealth[length=4pt]}, dashed, draw=medicaGreen] (note) -- node[midway, above, font=\tiny, text=medicaGreen] {\guillemotleft include\guillemotright} (calc);
  \draw[-{Stealth[length=4pt]}, dashed, draw=medicaGreen] (calc) -- node[midway, right, font=\tiny, text=medicaGreen] {\guillemotleft extend\guillemotright} (bloq);
\end{tikzpicture}
\caption{Cas d'utilisation du diagnostic et du calcul de readiness}
\label{fig:uc-diagnostic}
\end{figure}

\subsection{Sous-système Portefeuille et Valeur}

\begin{figure}[H]
\centering
\begin{tikzpicture}[scale=0.95, transform shape]
  \node[acteur] (dir) at (0,2.2) {Direction};
  \node[acteur] (fin) at (0,0.8) {Finance};
  \node[acteur] (pmo) at (0,-0.6) {PMO};

  \begin{scope}[on background layer]
    \draw[draw=medicaBorder, fill=medicaBg, rounded corners=8pt] (2.8,-2.2) rectangle (10.4,3.1);
  \end{scope}
  \node[font=\small\bfseries, text=medicaDeep] at (6.6,2.75) {Modules M7 / M8 / M9 / M10};

  \node[cas] (prio)  at (4.6,2.0) {Ajuster les pondérations};
  \node[cas] (recalc)at (8.4,2.0) {Recalculer le classement};
  \node[cas] (fiche) at (4.6,1.0) {Générer la fiche projet};
  \node[cas] (dec)   at (8.4,1.0) {Rendre la décision de comité};
  \node[cas] (cout)  at (4.6,0.0) {Saisir coûts et gains};
  \node[cas] (bc)    at (8.4,0.0) {Calculer le business case};
  \node[cas] (gate)  at (4.6,-1.0) {Passer une gate};
  \node[cas] (suivi) at (8.4,-1.0) {Suivre l'avancement};
  \node[cas] (roi)   at (6.5,-1.9) {Comparer valeur promise et réalisée};

  \draw[lien] (dir.east) -- (prio.west);
  \draw[lien] (dir.east) -- (fiche.west);
  \draw[lien] (fin.east) -- (cout.west);
  \draw[lien] (fin.east) -- (roi.west);
  \draw[lien] (pmo.east) -- (gate.west);
  \draw[lien] (pmo.east) -- (suivi.west);

  \draw[-{Stealth[length=4pt]}, dashed, draw=medicaGreen] (prio) -- node[midway, above, font=\tiny, text=medicaGreen] {\guillemotleft include\guillemotright} (recalc);
  \draw[-{Stealth[length=4pt]}, dashed, draw=medicaGreen] (cout) -- node[midway, above, font=\tiny, text=medicaGreen] {\guillemotleft include\guillemotright} (bc);
  \draw[-{Stealth[length=4pt]}, dashed, draw=medicaGreen] (fiche) -- node[midway, above, font=\tiny, text=medicaGreen] {\guillemotleft extend\guillemotright} (dec);
\end{tikzpicture}
\caption{Cas d'utilisation de la priorisation, de la fiche projet et de la valeur}
\label{fig:uc-portefeuille}
\end{figure}

\section{Modèle de données}

Le schéma compte \textbf{34~tables}, dont quatre tables d'association. Les diagrammes qui
suivent le découpent en quatre domaines pour rester lisibles. Les noms sont ceux des
classes SQLAlchemy réelles, la table physique étant précisée lorsqu'elle diffère.

\subsection{Domaine A — Organisation et contrôle d'accès}

\begin{figure}[H]
\centering
\begin{tikzpicture}[scale=0.9, transform shape, every node/.style={font=\scriptsize}]
  \node[classe] (org) at (0,0) {\umlclasse{Organization}{
    \underline{id} : int\\ name, code (unique)\\ industry, country\\ description}};
  \node[classe] (site) at (-4.6,-2.9) {\umlclasse{Site}{
    \underline{id} : int\\ \textit{organization\_id} FK\\ name, code\\ city, country, site\_type}};
  \node[classe] (dept) at (0,-2.9) {\umlclasse{Department}{
    \underline{id} : int\\ \textit{organization\_id} FK\\ \textit{site\_id} FK\\ name, code, family\\ head\_name}};
  \node[classe] (user) at (5,-2.9) {\umlclasse{User (app\_user)}{
    \underline{id} : int\\ \textit{organization\_id} FK\\ \textit{department\_id} FK\\ email (unique)\\ full\_name, job\_title\\ hashed\_password\\ is\_active, is\_superuser\\ last\_login\_at}};
  \node[classe] (role) at (5,-6.6) {\umlclasse{Role}{
    \underline{id} : int\\ \textit{organization\_id} FK\\ code, label\\ scope\_level\\ is\_system, is\_platform}};
  \node[classe] (perm) at (0,-6.6) {\umlclasse{Permission}{
    \underline{id} : int\\ code (unique)\\ label, module, group}};

  \draw[lien] (org.west) -| node[near start, above, font=\tiny] {1..*} (site.north);
  \draw[lien] (org.south) -- node[midway, right, font=\tiny] {1..*} (dept.north);
  \draw[lien] (org.east) -| node[near start, above, font=\tiny] {1..*} (user.north);
  \draw[lien] (org.east) -| ([xshift=0.6cm]role.north);
  \draw[lien] (dept.east) -- node[midway, above, font=\tiny] {0..1} (user.west);
  \draw[lien] (user.south) -- node[midway, right, font=\tiny] {\textit{user\_role} *..*} (role.north);
  \draw[lien] (role.west) -- node[midway, above, font=\tiny] {\textit{role\_permission} *..*} (perm.east);
  \draw[lien, dashed] (user.west) -- node[midway, below, font=\tiny] {\textit{user\_department} *..*} (dept.east);
\end{tikzpicture}
\caption{Domaine A — Organisation et contrôle d'accès}
\label{fig:classes-rbac}
\end{figure}

Quatre points méritent commentaire.

\begin{itemize}
  \item La table \code{User} s'appelle physiquement \code{app\_user}, \code{user} étant
    un mot réservé en PostgreSQL.
  \item Les permissions sont \emph{globales à la plateforme} : seuls les rôles, qui les
    regroupent, portent un \code{organization\_id}. Deux organisations partagent donc le
    même catalogue de permissions mais ont chacune leurs rôles.
  \item \code{User.department\_id} coexiste avec la table d'association
    \code{user\_department}. Le premier désigne le département principal et sert de
    repli ; la seconde porte les rattachements multiples. Cette coexistence est le
    résultat d'une migration qui devait préserver les données existantes.
  \item \code{Role.scope\_level} porte le périmètre par défaut. Le périmètre effectif
    d'un utilisateur est le plus large de ceux de ses rôles.
\end{itemize}

\subsection{Domaine B — Diagnostic, architecture et opportunités}

\begin{figure}[H]
\centering
\begin{tikzpicture}[scale=0.88, transform shape, every node/.style={font=\scriptsize}]
  \node[classe] (ass) at (0,0) {\umlclasse{Assessment}{
    \underline{id} : int\\ \textit{organization\_id} FK\\ \textit{site\_id}, \textit{department\_id} FK\\ name, campaign, status\\ blocking\_threshold\\ completion\_rate}};
  \node[classe] (resp) at (0,-3.4) {\umlclasse{AssessmentResponse}{
    \underline{id} : int\\ \textit{assessment\_id} FK\\ axis, dimension\_code\\ dimension\_label, question\\ score (0--5), weight\\ confidence (0--1)\\ evidence, is\_critical}};
  \node[classe] (opp) at (5.8,-1.2) {\umlclasse{Opportunity}{
    \underline{id} : int\\ \textit{organization\_id} FK\\ \textit{department\_id}, \textit{site\_id} FK\\ \textit{process\_id} FK\\ code, title, summary\\ as\_is\_context, pain\_points\\ answers : JSON\\ 15 faits de qualification\\ 6 composantes de score\\ risk\_probability, risk\_impact\\ priority\_score, rank, wave\\ score\_explanation : JSON}};
  \node[classe] (reco) at (11.4,0.2) {\umlclasse{AIRecommendation}{
    \underline{id} : int\\ \textit{opportunity\_id} FK (1:1)\\ ai\_type, ai\_type\_label\\ solution\_pattern\\ rationale, confidence\\ fired\_rules : JSON\\ alternatives : JSON\\ guardrails : JSON}};
  \node[classe] (dep) at (11.4,-3.4) {\umlclasse{Dependency}{
    \underline{id} : int\\ \textit{source\_opportunity\_id} FK\\ \textit{target\_opportunity\_id} FK\\ kind, asset\_ref, label\\ strength, is\_blocking}};
  \node[classe] (arch) at (5.8,-5.9) {\umlclasse{Process / Application / DataSource / KPI}{
    Référentiel AS-IS, chacun rattaché à l'organisation\\
    DataSource : quality\_score, has\_ground\_truth\\
    KPI : baseline, target, current}};

  \draw[lien] (ass.south) -- node[midway, right, font=\tiny] {1..*} (resp.north);
  \draw[lien] (opp.east) -- node[midway, above, font=\tiny] {1..1} (reco.west);
  \draw[lien] (opp.east) -| node[near start, below, font=\tiny] {*..*} (dep.west);
  \draw[lien] (opp.south) -- node[midway, right, font=\tiny] {0..1} (arch.north);
\end{tikzpicture}
\caption{Domaine B — Diagnostic, référentiel d'architecture et opportunités}
\label{fig:classes-diagnostic}
\end{figure}

La table \code{opportunity} est la plus large du schéma avec 48~colonnes. Elle regroupe
trois familles d'information : le contexte métier rédigé, quinze faits booléens ou
numériques qui alimentent le moteur de qualification, et six composantes de score
assorties du résultat calculé. Le champ \code{score\_explanation}, de type JSON,
conserve la décomposition du score au moment du dernier calcul.

\code{AIRecommendation} entretient une relation un à un avec l'opportunité. Son nom
contient « AI » parce qu'elle décrit un type de solution d'intelligence artificielle ;
elle est produite par un moteur de règles, non par un modèle.

\subsection{Domaine C — Projets et valeur}

\begin{figure}[H]
\centering
\begin{tikzpicture}[scale=0.88, transform shape, every node/.style={font=\scriptsize}]
  \node[classe] (proj) at (0,0) {\umlclasse{Project}{
    \underline{id} : int\\ \textit{organization\_id} FK\\ \textit{opportunity\_id} FK (1:1)\\ \textit{department\_id} FK\\ code, name, phase, status\\ sponsor, owners\\ budget\_eur, committed\_eur\\ progress, next\_gate\_date}};
  \node[classe] (stage) at (5.4,1.4) {\umlclasse{Stage}{
    \underline{id} : int\\ \textit{project\_id} FK\\ code, name, sequence\\ status, gate\_passed\\ gate\_criteria : JSON}};
  \node[classe] (ms) at (5.4,-1.0) {\umlclasse{Milestone}{
    \underline{id} : int\\ \textit{project\_id}, \textit{stage\_id} FK\\ name, due\_date\\ status, is\_gate}};
  \node[classe] (task) at (5.4,-3.2) {\umlclasse{Task}{
    \underline{id} : int\\ \textit{project\_id}, \textit{stage\_id} FK\\ title, assignee\\ status, due\_date, effort}};
  \node[classe] (risk) at (0,-3.6) {\umlclasse{Risk}{
    \underline{id} : int\\ \textit{project\_id} FK\\ \textit{opportunity\_id} FK\\ title, category\\ probability (1--5)\\ impact (1--5)\\ severity, status}};
  \node[classe] (mit) at (0,-6.4) {\umlclasse{Mitigation}{
    \underline{id} : int\\ \textit{risk\_id} FK\\ action, owner\\ due\_date, status\\ residual\_prob/impact}};
  \node[classe] (cost) at (10.6,0.6) {\umlclasse{Cost}{
    \underline{id} : int\\ \textit{project\_id} FK\\ \textit{opportunity\_id} FK\\ label, category, nature\\ amount\_eur, recurring}};
  \node[classe] (ben) at (10.6,-2.2) {\umlclasse{Benefit}{
    \underline{id} : int\\ \textit{opportunity\_id} FK\\ \textit{kpi\_id} FK\\ annual\_amount\_eur\\ realization\_rate\\ pnl\_line}};
  \node[classe] (meas) at (10.6,-5.0) {\umlclasse{Measurement}{
    \underline{id} : int\\ \textit{kpi\_id} FK\\ measured\_on, value\\ period\_label, source}};

  \draw[lien] (proj.east) -- node[midway, above, font=\tiny] {1..*} (stage.west);
  \draw[lien] (proj.east) -- node[midway, above, font=\tiny] {1..*} (ms.west);
  \draw[lien] (proj.east) -- node[midway, below, font=\tiny] {1..*} (task.west);
  \draw[lien] (proj.south) -- node[midway, right, font=\tiny] {1..*} (risk.north);
  \draw[lien] (risk.south) -- node[midway, right, font=\tiny] {1..*} (mit.north);
  \draw[lien] (stage.east) -| (cost.north);
  \draw[lien] (ben.west) -- node[midway, above, font=\tiny] {0..1} (task.east);
\end{tikzpicture}
\caption{Domaine C — Cycle de vie des projets et valorisation}
\label{fig:classes-projets}
\end{figure}

\code{Cost} et \code{Benefit} peuvent être rattachés soit à une opportunité, soit à un
projet. Dans le jeu de démonstration ils le sont à l'opportunité, ce qui permet de
calculer un business case avant même que le projet n'existe.

\subsection{Domaine D — Gouvernance documentaire et audit}

\begin{figure}[H]
\centering
\begin{tikzpicture}[scale=0.92, transform shape, every node/.style={font=\scriptsize}]
  \node[classe] (doc) at (0,0) {\umlclasse{Document}{
    \underline{id} : int\\ \textit{organization\_id} FK\\ \textit{opportunity\_id} FK\\ \textit{project\_id} FK\\ code, title, doc\_type\\ status, current\_version\\ author}};
  \node[classe] (ver) at (5.4,1.2) {\umlclasse{Version}{
    \underline{id} : int\\ \textit{document\_id} FK\\ number, label\\ content : JSON\\ change\_summary\\ created\_by, created\_at}};
  \node[classe] (app) at (5.4,-1.6) {\umlclasse{Approval}{
    \underline{id} : int\\ \textit{document\_id} FK\\ \textit{opportunity\_id} FK\\ \textit{project\_id} FK\\ \textit{approver\_user\_id} FK\\ stage\_code, committee\\ decision, decided\_on}};
  \node[classe] (log) at (0,-4.2) {\umlclasse{AuditLog — insertion seule}{
    \underline{id} : int\\ organization\_id, user\_id\\ user\_email\\ action, entity\_type\\ entity\_id, entity\_label\\ old\_value : JSON\\ new\_value : JSON\\ ip\_address, created\_at}};
  \node[classe] (notif) at (6,-4.6) {\umlclasse{Notification}{
    \underline{id} : int\\ \textit{organization\_id} FK\\ \textit{user\_id} FK\\ title, body, level\\ entity\_type, entity\_id\\ is\_read}};

  \draw[lien] (doc.east) -- node[midway, above, font=\tiny] {1..*} (ver.west);
  \draw[lien] (doc.east) -- node[midway, below, font=\tiny] {1..*} (app.west);
  \draw[lien, dashed, draw=medicaAmber] (log.north) -- node[midway, right, font=\tiny, text=medicaAmber]
    {trace toute écriture} (doc.south);
\end{tikzpicture}
\caption{Domaine D — Documents versionnés, approbations et journal d'audit}
\label{fig:classes-audit}
\end{figure}

\code{AuditLog} est volontairement dépourvue de clés étrangères. Une entrée d'audit doit
survivre à la suppression de l'entité qu'elle décrit : si l'opportunité est effacée, la
trace de sa suppression doit subsister. Les identifiants y sont donc de simples entiers,
sans contrainte référentielle.

\section{Architecture technique}

\subsection{Vue d'ensemble}

\begin{figure}[H]
\centering
\begin{tikzpicture}[node distance=0.65cm, scale=0.95, transform shape]
  \node[blocsombre, text width=3.4cm] (user) {\textbf{Utilisateur}\\\scriptsize navigateur};

  \node[bloc, below=of user, text width=10.5cm, fill=medicaSoft, draw=medicaGreen] (front)
    {\textbf{Frontend React 18 + TypeScript} \\[2pt]
     \scriptsize Vite · React Router · TanStack Query · Recharts · axios};

  \node[bloc, below=of front, text width=10.5cm, fill=white, draw=medicaTeal] (api)
    {\textbf{API REST — 99 endpoints sous \code{/api/v1}} \\[2pt]
     \scriptsize JSON · OpenAPI sur \code{/docs} · CORS restreint aux origines déclarées};

  \node[bloc, below=of api, text width=10.5cm, fill=white, draw=medicaGreen] (fastapi)
    {\textbf{FastAPI} — routers par module (M1 à M12) \\[2pt]
     \scriptsize Dépendances : session, utilisateur courant, permission, périmètre};

  \node[bloc, below=of fastapi, text width=10.5cm, fill=medicaSoft, draw=medicaGreen] (services)
    {\textbf{Services métier} — sans dépendance à FastAPI \\[2pt]
     \scriptsize readiness · recommendation · scoring · crossing · value · questionnaire · audit};

  \node[bloc, below=of services, text width=10.5cm, fill=white, draw=medicaTeal] (orm)
    {\textbf{SQLAlchemy 2.0} — 34 tables, sessions par requête \\[2pt]
     \scriptsize Migrations versionnées par Alembic};

  \node[blocsombre, below=of orm, text width=10.5cm] (db)
    {\textbf{PostgreSQL} \\[2pt]\scriptsize Filtrage \code{organization\_id} sur toute table métier};

  \draw[flechevert] (user) -- (front);
  \draw[flechevert] (front) -- node[right, font=\scriptsize, text=medicaMuted] {HTTPS + jeton JWT} (api);
  \draw[flechevert] (api) -- (fastapi);
  \draw[flechevert] (fastapi) -- (services);
  \draw[flechevert] (services) -- (orm);
  \draw[flechevert] (orm) -- (db);

  % éléments transverses
  \node[bloc, right=1.1cm of fastapi, text width=2.6cm, fill=medicaBg, draw=medicaAmber, font=\scriptsize]
    (sec) {\textbf{Sécurité}\\ JWT HS256\\ bcrypt\\ RBAC\\ périmètre};
  \draw[fleche, dashed, draw=medicaAmber] (sec) -- (fastapi);

  \node[bloc, right=1.1cm of db, text width=2.6cm, fill=medicaBg, draw=medicaTeal, font=\scriptsize]
    (docker) {\textbf{Docker}\\ postgres:16\\ backend\\ frontend};
  \draw[fleche, dashed, draw=medicaTeal] (docker) -- (db);
\end{tikzpicture}
\caption{Architecture en couches de la plateforme}
\label{fig:architecture}
\end{figure}

\subsection{Frontend}

L'interface est une application monopage construite avec React~18 et TypeScript,
servie par Vite. Le routage est assuré par React Router, les appels réseau par axios
encapsulé dans des hooks TanStack Query, et les graphiques par Recharts.

Le point structurant est que \emph{la navigation dérive des permissions}. Le composant
\code{Layout} filtre les entrées de menu selon les permissions renvoyées par
\code{/auth/me} ; une entrée n'apparaît que si son écran est réellement accessible. Le
même principe gouverne les boutons d'action.

\subsection{Backend}

L'API est écrite en Python~3.12 avec FastAPI. Elle est organisée en dix routers, un par
domaine fonctionnel, montés sous le préfixe \code{/api/v1}.

La logique métier vit dans sept services qui n'importent rien de FastAPI. Cette
séparation a une conséquence pratique : les formules de readiness, de score et de ROI
se testent en appelant directement les fonctions, sans serveur HTTP ni base de données.

\subsection{Base de données}

PostgreSQL héberge les 34~tables. Le schéma est versionné par trois migrations Alembic,
écrites pour s'appliquer sur une base déjà peuplée. La troisième, qui a introduit les
périmètres, crée les colonnes avec une valeur par défaut serveur le temps de la reprise,
recopie les rattachements existants depuis \code{department\_id}, puis retire la valeur
par défaut.

\subsection{API}

L'API compte 99~endpoints répartis comme suit.

\begin{table}[H]
\centering
\caption{Répartition des endpoints par domaine}
\label{tab:endpoints}
\small
\begin{tabular}{@{}L{3.6cm} C{1.6cm} L{8.2cm}@{}}
\toprule
\textbf{Préfixe} & \textbf{Nombre} & \textbf{Objet} \\
\midrule
\code{/org}           & 18 & Organisation, sites, départements, utilisateurs, rôles \\
\code{/projects}      & 16 & Projets, phases, jalons, tâches, risques, gates \\
\code{/crossings}     & 14 & Processus, applications, données, dépendances, doublons \\
\code{/opportunities} & 12 & Opportunités, qualification, priorisation \\
\code{/value}         &  9 & Coûts, bénéfices, mesures, business case \\
\code{/diagnostic}    &  8 & Campagnes, réponses, readiness \\
\code{/sheets}        &  7 & Fiches projet, versions, approbations, export \\
\code{/auth}          &  5 & Connexion, profil, mot de passe \\
\code{/admin}         &  4 & Journal d'audit, notifications \\
\code{/dashboards}    &  4 & Agrégations Direction, département, portefeuille \\
Système               &  2 & Sonde de santé et racine \\
\midrule
\textbf{Total}        & \textbf{99} & \\
\bottomrule
\end{tabular}
\end{table}

\subsection{Déploiement}

Un fichier \code{docker-compose.yml} décrit trois services : PostgreSQL~16 avec sonde
de santé et volume persistant, le backend construit depuis \code{backend/Dockerfile}, et
le frontend depuis \code{frontend/Dockerfile}. Le backend attend que la base soit saine
avant de démarrer.

\begin{remarque}[Sur la version de PostgreSQL]
Le fichier Docker déclare \code{postgres:16-alpine}. L'instance utilisée pendant le
développement est une PostgreSQL~12.4 installée localement. Aucune fonctionnalité
spécifique à la version~16 n'est employée, mais cet écart mérite d'être connu avant une
mise en production.
\end{remarque}

% =====================================================================
%  Chapitre 3 — Réalisation
% =====================================================================
\chapter{Réalisation}
\label{chap:realisation}

\section{Environnement et technologies}

Les versions ci-dessous sont celles qui figurent dans \code{backend/requirements.txt} et
\code{frontend/package.json}. Lorsqu'une version n'est pas épinglée dans un fichier du
projet, la mention correspondante l'indique.

\begin{longtable}{@{}L{4.2cm} L{2.9cm} L{7.1cm}@{}}
\caption{Technologies employées et versions réellement déclarées}
\label{tab:technos}\\
\toprule
\textbf{Technologie} & \textbf{Version} & \textbf{Rôle dans le projet} \\
\midrule
\endfirsthead
\multicolumn{3}{c}{\small\itshape suite du tableau~\ref{tab:technos}}\\
\toprule
\textbf{Technologie} & \textbf{Version} & \textbf{Rôle} \\
\midrule
\endhead
\bottomrule
\endfoot

\multicolumn{3}{@{}l}{\textit{\textbf{Backend}}}\\
Python      & 3.12.7 & Langage du serveur (version de l'interpréteur utilisé) \\
FastAPI     & 0.115.6 & Framework web, validation des entrées, OpenAPI \\
Uvicorn     & 0.34.0 & Serveur ASGI \\
SQLAlchemy  & 2.0.36 & ORM, style déclaratif typé \\
Alembic     & 1.14.0 & Migrations de schéma versionnées \\
psycopg2-binary & 2.9.10 & Pilote PostgreSQL \\
Pydantic    & 2.10.4 & Schémas d'entrée et de sortie \\
pydantic-settings & 2.7.0 & Configuration par variables d'environnement \\
python-jose & 3.3.0 & Signature et vérification des jetons JWT \\
passlib     & 1.7.4 & Interface de hachage \\
bcrypt      & 4.0.1 & Algorithme de hachage des mots de passe \\
email-validator & 2.2.0 & Validation des adresses électroniques \\
python-multipart & 0.0.20 & Formulaire du flux OAuth2 \\
pytest      & 8.3.4 & Exécution des tests \\
httpx       & 0.28.1 & Client HTTP de test \\
\midrule
\multicolumn{3}{@{}l}{\textit{\textbf{Frontend}}}\\
Node.js     & 24.12.0 & Environnement d'exécution (version utilisée) \\
npm         & 11.6.2 & Gestionnaire de paquets \\
React       & 18.3 & Bibliothèque d'interface \\
TypeScript  & 5.7 & Typage statique \\
Vite        & 6.0 & Serveur de développement et empaqueteur \\
React Router & 6.28 & Routage côté client \\
TanStack Query & 5.62 & Cache et synchronisation des données serveur \\
Recharts    & 2.13 & Graphiques \\
axios       & 1.7 & Client HTTP \\
\midrule
\multicolumn{3}{@{}l}{\textit{\textbf{Infrastructure}}}\\
PostgreSQL  & 16-alpine (Docker) & Base déclarée dans \code{docker-compose.yml} \\
PostgreSQL  & 12.4 (local) & Instance effectivement utilisée en développement \\
Docker Compose & non versionné & Orchestration des trois services \\
\end{longtable}

\begin{remarque}[Lecture des versions du frontend]
Les dépendances JavaScript sont déclarées avec l'opérateur \code{\textasciicircum{}},
qui autorise les montées de version mineures. Le tableau retient la version minimale
déclarée. La version exactement installée figure dans \code{package-lock.json}.
\end{remarque}

\section{Organisation du code}

\begin{figure}[H]
\centering
\begin{minipage}{0.92\textwidth}
\small\ttfamily\raggedright
\textcolor{medicaDeep}{backend/app/}\\
\hspace*{1em}core/ \hfill \textcolor{medicaMuted}{\rmfamily config, security (JWT, bcrypt), rbac (catalogue)}\\
\hspace*{1em}db/ \hfill \textcolor{medicaMuted}{\rmfamily session, base déclarative, seed}\\
\hspace*{1em}models/ \hfill \textcolor{medicaMuted}{\rmfamily 34 tables réparties en 7 modules}\\
\hspace*{1em}schemas/ \hfill \textcolor{medicaMuted}{\rmfamily schémas Pydantic d'entrée et de sortie}\\
\hspace*{1em}services/ \hfill \textcolor{medicaMuted}{\rmfamily 7 services métier, sans dépendance HTTP}\\
\hspace*{1em}api/deps.py \hfill \textcolor{medicaMuted}{\rmfamily session, utilisateur, permission, périmètre}\\
\hspace*{1em}api/routers/ \hfill \textcolor{medicaMuted}{\rmfamily 10 routers, 99 endpoints}\\
\hspace*{1em}main.py \hfill \textcolor{medicaMuted}{\rmfamily application, CORS, montage des routers}\\[4pt]
\textcolor{medicaDeep}{backend/alembic/versions/} \hfill \textcolor{medicaMuted}{\rmfamily 3 migrations}\\
\textcolor{medicaDeep}{backend/tests/} \hfill \textcolor{medicaMuted}{\rmfamily 7 fichiers, 112 tests}\\[6pt]
\textcolor{medicaDeep}{frontend/src/}\\
\hspace*{1em}lib/ \hfill \textcolor{medicaMuted}{\rmfamily client axios, contexte d'authentification, hooks}\\
\hspace*{1em}components/ \hfill \textcolor{medicaMuted}{\rmfamily mise en page, graphiques, panneaux, icônes}\\
\hspace*{1em}pages/ \hfill \textcolor{medicaMuted}{\rmfamily un composant par écran}\\
\hspace*{1em}styles.css \hfill \textcolor{medicaMuted}{\rmfamily charte graphique, variables CSS}
\end{minipage}
\caption{Arborescence des répertoires significatifs}
\label{fig:arborescence}
\end{figure}

La règle d'organisation du backend est simple : un router ne calcule rien. Il vérifie
les droits, charge les entités, délègue à un service, journalise et renvoie. Toute
formule est dans \code{services/}.

% =====================================================================
\section{L'intelligence de la plateforme : un moteur de règles, pas un modèle}
\label{sec:pas-dia}

Cette section est la plus importante du rapport en matière d'honnêteté technique. Le
produit s'appelle \emph{AOI Transformation}, une table de la base s'appelle
\code{ai\_recommendation}, et la plateforme recommande des types de solutions
d'intelligence artificielle. Rien de tout cela n'implique qu'elle en contienne une.

\subsection{Le constat, établi sur le code}

Une recherche exhaustive a été menée sur l'ensemble du dépôt, portant sur les
bibliothèques d'apprentissage automatique (\code{scikit-learn}, \code{tensorflow},
\code{torch}, \code{transformers}, \code{keras}, \code{xgboost}, \code{pandas},
\code{numpy}), sur les clients de fournisseurs de modèles (\code{openai},
\code{anthropic}, \code{mistralai}, \code{cohere}, \code{huggingface}), ainsi que sur
tout appel réseau sortant du backend. Le résultat est sans ambiguïté :

\begin{itemize}
  \item aucune de ces bibliothèques n'apparaît dans \code{requirements.txt} ni dans
    \code{package.json} ;
  \item aucun import correspondant n'existe dans le code ;
  \item aucun fichier de modèle entraîné (\code{.pkl}, \code{.h5}, \code{.onnx},
    \code{.pt}) n'est présent dans le dépôt ;
  \item aucune clé d'API de fournisseur de modèle n'est attendue par la configuration ;
  \item le backend n'émet aucune requête HTTP sortante. Le seul client HTTP installé,
    \code{httpx}, sert exclusivement aux tests.
\end{itemize}

Les deux seules occurrences du sigle « LLM » dans le code sont des chaînes de caractères
en français, décrivant le \emph{patron de solution} recommandé à l'utilisateur —
« Assistant LLM ancré sur base documentaire (RAG), sortie validée ». Ce sont des
recommandations adressées à un humain, pas des appels à un modèle.

\begin{remarque}[Formulation exacte à retenir]
La plateforme \emph{recommande} des solutions d'intelligence artificielle. Elle n'en
\emph{utilise} aucune. Son moteur de qualification est un système à base de règles
pondérées, entièrement déterministe, dont chaque sortie est reconstituable à la main.
\end{remarque}

\subsection{Moteur de règles et apprentissage automatique : la distinction}

La confusion entre les deux est fréquente, y compris dans la littérature commerciale.
Elle porte sur l'origine de la connaissance.

\begin{table}[H]
\centering
\caption{Ce qui sépare un moteur de règles d'un modèle d'apprentissage}
\label{tab:regles-vs-ml}
\small
\begin{tabular}{@{}L{3.3cm} L{5.3cm} L{5.3cm}@{}}
\toprule
& \textbf{Moteur de règles (ce projet)} & \textbf{Apprentissage automatique} \\
\midrule
Origine de la connaissance
& Écrite par un expert humain, lisible dans le code
& Inférée statistiquement depuis des données d'entraînement \\
Forme
& Prédicats booléens pondérés
& Paramètres numériques appris, souvent des millions \\
Besoin en données
& Aucun jeu d'entraînement
& Corpus annoté, représentatif et volumineux \\
Déterminisme
& Deux exécutions sur les mêmes faits donnent le même résultat
& Résultat dépendant de l'entraînement et du tirage \\
Explicabilité
& Native : la règle qui a décidé est nommée
& À reconstruire après coup (importance des variables, SHAP, LIME) \\
Évolution
& Modifier une règle, une ligne de code, relire la revue
& Réentraîner, réévaluer, surveiller la dérive \\
Validation réglementaire
& Le jeu de règles est un document versionnable et signable
& Exige un dossier de validation et des jeux de test figés \\
Limite principale
& Ne découvre rien qu'un expert n'ait prévu
& Capte des régularités qu'aucun expert n'aurait formulées \\
\bottomrule
\end{tabular}
\end{table}

Le moteur de ce projet relève de la première colonne. On peut le qualifier de système
expert au sens classique : une base de connaissances — 39~règles réparties sur 8~types
de solution — et un mécanisme d'inférence — l'agrégation pondérée et l'application des
règles bloquantes.

\subsection{Pourquoi ce choix est adapté ici}

Ce n'est pas un pis-aller. Quatre arguments le soutiennent.

\textbf{Il n'y a pas de données d'entraînement.} Entraîner un modèle à prédire
« quel type de solution convient à ce besoin » supposerait un historique de plusieurs
centaines de besoins qualifiés puis validés par leur résultat réel. Aucune organisation
ne dispose de ce corpus au moment où elle démarre sa démarche d'adoption. Le problème
est, par construction, un problème de démarrage à froid.

\textbf{La décision doit être opposable.} La recommandation est présentée en comité
d'investissement. Un expert du domaine doit pouvoir contester une règle précise. Avec
un modèle statistique, la discussion se déplace vers la qualité du jeu d'entraînement,
terrain sur lequel un comité métier ne peut pas arbitrer.

\textbf{Le contexte réglementaire pénalise l'opacité.} En périmètre GxP, une décision
non explicable est une décision à documenter lourdement. Un jeu de règles versionné est
un artefact de validation naturel.

\textbf{Le déterminisme rend le système testable.} Les règles sont couvertes par des
tests qui vérifient des sorties exactes. Un modèle probabiliste ne se teste pas ainsi ;
il se mesure sur des métriques agrégées.

\subsection{Ce que cela implique, et l'évolution envisageable}

La contrepartie est réelle : le moteur ne découvre rien. Il ne peut pas proposer un type
de solution auquel les règles n'ont pas pensé, ni apprendre des décisions passées. Si
tous les comités écartent systématiquement les recommandations « agentique », le moteur
continuera à les produire.

L'intégration d'un modèle de langage est envisageable et serait utile sur trois points
précis : pré-remplir les faits de qualification à partir d'une description libre du
besoin, rédiger un premier jet de fiche projet, et détecter des doublons sémantiques
entre opportunités formulées différemment. Ces trois usages sont décrits au
chapitre~\ref{chap:conclusion} comme \emph{perspectives}. Aucun n'est implémenté, et
aucun ne devrait remplacer la décision du moteur de règles : le modèle proposerait des
entrées, les règles continueraient de décider.

% =====================================================================
\section{Le moteur de qualification (M4)}

\subsection{Principe}

Le moteur compare une opportunité à huit types de solution. Il ne choisit pas « le plus
probable » ; il évalue chaque type indépendamment, puis retient celui dont le plus grand
poids de règles est satisfait.

Les huit types sont : IA classique (apprentissage supervisé, séries temporelles), vision
par ordinateur, traitement du langage, IA générative, IA agentique, décisionnel,
automatisation robotisée de processus, et refonte de processus. Ce dernier mérite d'être
souligné : c'est la réponse « pas d'IA », et le moteur la produit lorsque le processus
est instable ou les données absentes.

\subsection{Les faits d'entrée}

Quinze faits sont extraits de l'opportunité. Treize sont booléens
(\code{data\_available}, \code{data\_structured}, \code{ground\_truth\_available},
\code{unstructured\_text}, \code{image\_or\_video}, \code{content\_generation},
\code{multi\_step\_autonomy}, \code{rule\_based}, \code{repetitive\_task},
\code{reporting\_need}, \code{process\_unstable}, \code{regulated\_gxp},
\code{human\_in\_the\_loop}), deux sont numériques (\code{data\_quality} sur~5 et
\code{volume\_per\_month}). Les réponses libres du questionnaire peuvent surcharger un
fait, ce qui permet d'affiner sans modifier le schéma.

\subsection{Le calcul}

Soit $R_t$ l'ensemble des règles du type~$t$, $w_r$ le poids de la règle~$r$ et
$\sigma_r \in \{0,1\}$ l'indicateur de satisfaction. Le score brut du type vaut :

\begin{equation}
s_t = \frac{\displaystyle\sum_{r \in R_t} w_r \, \sigma_r}{\displaystyle\sum_{r \in R_t} w_r}
\label{eq:score-type}
\end{equation}

Une règle peut être déclarée \emph{bloquante}. Si une règle bloquante n'est pas
satisfaite, le type est disqualifié et son score divisé par quatre :

\begin{equation}
s_t^{\,\text{final}} =
\begin{cases}
s_t & \text{si toutes les règles bloquantes de } t \text{ sont satisfaites,}\\[6pt]
0{,}25 \times s_t & \text{sinon.}
\end{cases}
\label{eq:disqualification}
\end{equation}

Le choix de diviser plutôt que d'annuler est délibéré. Un type disqualifié reste visible
dans les alternatives avec son score pénalisé et le motif de sa disqualification. Une
opportunité de vision par ordinateur sans images annotées apparaît donc comme
« pertinente mais bloquée sur la vérité terrain », information bien plus utile qu'une
absence.

Le type retenu est celui de score final maximal ; en cas d'égalité, l'ordre alphabétique
du code technique départage, ce qui garantit un résultat reproductible. Les trois
suivants sont renvoyés comme alternatives.

\subsection{Enchaînement}

\begin{figure}[H]
\centering
\begin{tikzpicture}[node distance=0.55cm, scale=0.95, transform shape]
  \node[bloc, text width=4.6cm, fill=medicaSoft, draw=medicaGreen] (opp)
    {Opportunité\\\scriptsize 15 faits de qualification};
  \node[bloc, below=of opp, text width=4.6cm] (facts)
    {Extraction des faits\\\scriptsize surcharge par le questionnaire};
  \node[bloc, below=of facts, text width=6.2cm, draw=medicaTeal] (loop)
    {Pour chacun des 8 types :\\\scriptsize évaluer les règles, cumuler les poids satisfaits};
  \node[bloc, below=of loop, text width=6.2cm] (bloq)
    {Règle bloquante non satisfaite ?\\\scriptsize si oui : score $\times$ 0,25 et motif conservé};
  \node[bloc, below=of bloq, text width=6.2cm] (tri)
    {Tri décroissant des 8 scores};
  \node[blocvert, below=of tri, text width=7cm] (out)
    {\textbf{Recommandation}\\\scriptsize type retenu · confiance · justification rédigée\\
     \scriptsize règles satisfaites et manquées · 3 alternatives · garde-fous};
  \node[bloc, right=1.2cm of out, text width=3cm, draw=medicaAmber, fill=medicaBg, font=\scriptsize]
    (gxp) {Si \code{regulated\_gxp} : 3 garde-fous GxP ajoutés en tête};

  \draw[flechevert] (opp) -- (facts);
  \draw[flechevert] (facts) -- (loop);
  \draw[flechevert] (loop) -- (bloq);
  \draw[flechevert] (bloq) -- (tri);
  \draw[flechevert] (tri) -- (out);
  \draw[fleche, dashed, draw=medicaAmber] (gxp) -- (out);
\end{tikzpicture}
\caption{Enchaînement du moteur de qualification}
\label{fig:moteur-qualification}
\end{figure}

\subsection{Le cœur de l'algorithme}

\begin{lstlisting}[language=Python, caption={Évaluation d'un type de solution, extrait de services/recommendation.py}, label={lst:evaluate}]
def evaluate_type(ai_type, facts):
    rules = RULES[ai_type]
    total_weight = sum(r.weight for r in rules)
    fired_weight = 0.0
    disqualified_by = None

    for rule in rules:
        satisfied = bool(rule.predicate(facts))
        if satisfied:
            fired_weight += rule.weight
            fired.append(...)          # trace conservee pour l'explication
        else:
            missed.append(...)
            if rule.blocking and disqualified_by is None:
                disqualified_by = rule.label

    score = fired_weight / total_weight
    if disqualified_by:
        score *= 0.25                  # le type reste visible, fortement penalise
    return TypeEvaluation(ai_type, score, fired, missed, disqualified_by)
\end{lstlisting}

L'élément notable n'est pas le calcul, qui tient en dix lignes, mais le fait que les
listes \code{fired} et \code{missed} soient conservées. C'est ce qui permet à l'API de
répondre « voici pourquoi » et non seulement « voici quoi ». Ces listes sont persistées
en JSON dans \code{ai\_recommendation.fired\_rules}, ce qui rend la recommandation
reconstituable des mois plus tard, même si les règles ont changé entre-temps.

\subsection{Illustration}

Une opportunité d'inspection visuelle de blisters déclare des images disponibles et
annotées, une tâche répétitive, un volume mensuel de plusieurs milliers d'unités et un
périmètre GxP. Les cinq règles du type \emph{vision par ordinateur} sont satisfaites, le
score vaut 1,00. Le type \emph{IA classique} est disqualifié par la règle CLS-2, les
données n'étant pas structurées ; il figure en alternative avec un score divisé par
quatre et le motif « Données structurées (ERP, MES, LIMS) ». Le périmètre GxP ajoute en
tête des garde-fous la validation informatique, la traçabilité ALCOA+ et la revue
humaine obligatoire.

% =====================================================================
\section{L'algorithme de readiness (M5)}

\subsection{Formules}

Chaque dimension du diagnostic porte une note $n_i \in [0,5]$, un poids $p_i$ et une
confiance $c_i \in [0,1]$ déclarée par le répondant. Le score global est :

\begin{equation}
\text{readiness} = 100 \times
\frac{\displaystyle\sum_{i} n_i \, p_i \, c_i}{\displaystyle\sum_{i} 5 \, p_i \, c_i}
\label{eq:readiness}
\end{equation}

La confiance apparaît au numérateur \emph{et} au dénominateur. Une réponse mal étayée
pèse donc moins dans les deux termes : elle ne fait ni monter ni descendre
artificiellement le score, elle s'efface.

La moyenne d'un axe — technique ou métier — n'utilise que le poids :

\begin{equation}
\bar{n}_{\text{axe}} =
\frac{\displaystyle\sum_{i \in \text{axe}} n_i \, p_i}{\displaystyle\sum_{i \in \text{axe}} p_i}
\label{eq:moyenne-axe}
\end{equation}

La différence de traitement est intentionnelle. Le score global mesure une maturité
\emph{consolidée}, pour laquelle la solidité du constat compte. La moyenne d'axe décrit
le niveau constaté, indépendamment de la qualité de l'instruction.

\subsection{Dimensions non instruites}

Une confiance nulle signale une dimension non encore évaluée. Elle n'apporte ni au
numérateur ni au dénominateur — une dimension non mesurée ne vaut pas zéro — et elle
pèse uniquement sur le taux de complétion :

\begin{equation}
\text{complétion} = 100 \times
\frac{\text{dimensions instruites}}{\text{dimensions du référentiel}}
\end{equation}

Ce taux est toujours calculé, jamais saisi. La conséquence pratique est importante :
ajouter une dimension au référentiel fait baisser le taux de complétion de toutes les
campagnes, ce qui est exact, sans faire baisser mécaniquement leur score, ce qui serait
faux.

\subsection{Dimensions bloquantes et points de vigilance}

Une dimension instruite dont la note passe sous le seuil de la campagne — 2,5 par
défaut — est signalée. Sa qualification dépend de son caractère critique :

\begin{equation}
\text{bloquante} \iff (n_i < \theta) \wedge \text{critique}_i
\qquad
\text{vigilance} \iff (n_i < \theta) \wedge \neg\,\text{critique}_i
\end{equation}

Une dimension bloquante interdit le passage en production tant qu'elle n'est pas
remédiée. Un point de vigilance relève du plan de conduite du changement. Le niveau de
maturité affiché est $\lfloor \text{readiness} / 20 \rfloor$, plafonné à~5, sur une
échelle allant de « Inexistant » à « Optimisé ».

% =====================================================================
\section{La priorisation du portefeuille (M7)}

\subsection{Le score de priorité}

Six critères notés sur 100 sont combinés, puis une pénalité de risque est retranchée :

\begin{equation}
\text{score} = \sum_{k \in K} \omega_k \, v_k \; - \; \pi
\qquad\text{avec}\qquad
\sum_{k \in K} \omega_k = 1
\label{eq:priorite}
\end{equation}

Les pondérations par défaut sont données au tableau~\ref{tab:poids}.

\begin{table}[H]
\centering
\caption{Pondérations par défaut des six critères de priorisation}
\label{tab:poids}
\small
\begin{tabular}{@{}L{5cm} C{2cm} L{6.4cm}@{}}
\toprule
\textbf{Critère} & \textbf{Poids} & \textbf{Ce qu'il mesure} \\
\midrule
Impact métier             & 0,25 & Gain opérationnel attendu sur le processus \\
Impact financier / P\&L   & 0,20 & Effet sur le compte de résultat \\
Faisabilité technique     & 0,20 & Disponibilité des données et des compétences \\
Alignement stratégique    & 0,15 & Cohérence avec les priorités de l'organisation \\
Readiness organisationnel & 0,10 & Maturité du département porteur \\
Synergies                 & 0,10 & Réutilisation d'actifs partagés \\
\midrule
\textbf{Total}            & \textbf{1,00} & \\
\bottomrule
\end{tabular}
\end{table}

Ces poids sont ajustables depuis l'interface. Toute pondération fournie est renormalisée
à~1 avant usage, ce qui garantit que les scores restent comparables quelle que soit la
saisie. Un utilisateur qui porte l'impact financier de 0,20 à 0,60 sans toucher au reste
obtient bien une grille valide, dont les autres poids ont été proportionnellement
réduits.

\subsection{La pénalité de risque}

Le risque brut est le produit de la probabilité et de l'impact, tous deux notés de 1
à~5, soit une valeur de 1 à~25. Cette valeur est ramenée à l'échelle du score :

\begin{equation}
\pi = \frac{P \times I}{25} \times \Pi_{\max},
\qquad \Pi_{\max} = 15
\label{eq:penalite}
\end{equation}

Une opportunité au risque maximal perd donc 15~points sur 100. Le plafond est un choix
d'étalonnage : assez élevé pour déclasser une opportunité risquée, assez bas pour qu'un
risque maîtrisable n'annule pas un impact fort. Le score final est borné dans $[0,100]$.

\subsection{Rangs, vagues et quadrants}

Les opportunités sont triées par score décroissant, l'identifiant départageant les
égalités. Le rang est la position dans ce classement ; la vague est
$\lfloor (\text{rang}-1) / 6 \rfloor + 1$, la taille de vague étant paramétrable. Ce
découpage donne une lecture de séquencement plutôt qu'un simple palmarès : la vague~1
est ce que l'organisation lance maintenant, la vague~2 ce qu'elle prépare.

La matrice de portefeuille croise impact métier et faisabilité autour d'un seuil de~50,
définissant quatre quadrants : \emph{prioritaire} (fort impact, faisable),
\emph{quick win} (impact modéré, faisable), \emph{à réfléchir} (fort impact, difficile)
et \emph{à écarter}. La taille du point représente le bénéfice annuel attendu.

% =====================================================================
\section{Le pilotage de la valeur (M10)}

\subsection{Coût total de possession}

Les coûts sont séparés en non récurrents et récurrents. Sur un horizon de $H$ années :

\begin{equation}
\text{TCO} = \underbrace{\sum_{c \,\in\, \text{non récurrents}} m_c}_{\text{CAPEX}}
\; + \;
\underbrace{\Big(\sum_{c \,\in\, \text{récurrents}} m_c\Big)}_{\text{OPEX annuel}} \times H
\label{eq:tco}
\end{equation}

\subsection{Gains et retour sur investissement}

Chaque bénéfice porte un montant annuel et un taux de réalisation, qui traduit la
prudence de l'hypothèse. Le gain annuel retenu est pondéré par ce taux :

\begin{equation}
G = \sum_{b} m_b \, \rho_b
\qquad\text{puis}\qquad
\text{ROI}\,(\%) = \frac{G \times H - \text{TCO}}{\text{TCO}} \times 100
\label{eq:roi}
\end{equation}

\subsection{Délai de retour}

\begin{equation}
\text{payback (mois)} = \frac{\text{CAPEX}}{\big(G - \text{OPEX annuel}\big) / 12}
\label{eq:payback}
\end{equation}

Le numérateur est le seul CAPEX, pas le TCO. Le payback répond à la question « en
combien de mois l'investissement initial est-il récupéré par le gain annuel net
d'exploitation ? ». Y faire figurer trois années d'OPEX qui ne sont pas encore engagées
allongerait artificiellement le délai. Le calcul n'est produit que si le gain net annuel
et le CAPEX sont strictement positifs ; sinon le délai est déclaré non calculable plutôt
qu'affiché à zéro.

\subsection{Valeur promise et valeur réalisée}

Les bénéfices sont rattachés à des indicateurs, dont les mesures successives sont
horodatées. Le taux d'atteinte d'un indicateur est
$(\text{courant} - \text{référence}) / (\text{cible} - \text{référence}) \times 100$,
borné à 150~\% pour absorber les dépassements sans fausser les agrégats. C'est cette
comparaison qui ferme la boucle du parcours : le business case cesse d'être un document
d'obtention de budget pour devenir un engagement suivi.

% =====================================================================
\section{Sécurité et isolation}

\subsection{Chaîne de contrôle d'une requête}

\begin{figure}[H]
\centering
\begin{tikzpicture}[node distance=0.5cm, scale=0.95, transform shape]
  \node[bloc, text width=6cm, fill=medicaSoft, draw=medicaGreen] (req)
    {Requête HTTP\\\scriptsize en-tête \code{Authorization: Bearer}};
  \node[bloc, below=of req, text width=6cm, draw=medicaTeal] (jwt)
    {Vérification du jeton\\\scriptsize signature HS256, expiration};
  \node[bloc, below=of jwt, text width=6cm, draw=medicaTeal] (user)
    {Chargement de l'utilisateur\\\scriptsize compte actif ? sinon 401};
  \node[bloc, below=of user, text width=6cm, draw=medicaTeal] (perm)
    {Permission requise par l'endpoint\\\scriptsize absente $\rightarrow$ \textbf{403}};
  \node[bloc, below=of perm, text width=6cm, draw=medicaTeal] (tenant)
    {Filtre \code{organization\_id} issu du jeton\\\scriptsize jamais de la requête};
  \node[bloc, below=of tenant, text width=6cm, draw=medicaTeal] (scope)
    {Périmètre de l'utilisateur\\\scriptsize lecture hors périmètre $\rightarrow$ \textbf{404}\\
     \scriptsize écriture hors périmètre $\rightarrow$ \textbf{403}};
  \node[blocvert, below=of scope, text width=6cm] (exec)
    {Exécution et journalisation\\\scriptsize entrée \code{audit\_log} si écriture};

  \foreach \a/\b in {req/jwt, jwt/user, user/perm, perm/tenant, tenant/scope, scope/exec}
    \draw[flechevert] (\a) -- (\b);

  \node[bloc, right=1.2cm of perm, text width=3.1cm, fill=medicaBg, draw=medicaRed, font=\scriptsize]
    (refus) {Refus\\ 401 · 403 · 404\\ aucune donnée renvoyée};
  \draw[fleche, draw=medicaRed] (perm.east) -- (refus.west);
  \draw[fleche, draw=medicaRed] (scope.east) -| (refus.south);
  \draw[fleche, draw=medicaRed] (jwt.east) -| (refus.north);
\end{tikzpicture}
\caption{Contrôles traversés par toute requête authentifiée}
\label{fig:securite}
\end{figure}

\subsection{Authentification}

Les mots de passe sont hachés avec bcrypt ; aucun mot de passe en clair n'est stocké ni
journalisé. La connexion émet un jeton JWT signé en HS256, dont la charge utile contient
l'identifiant de l'utilisateur, celui de son organisation, son rôle principal, la date
d'émission et la date d'expiration. La durée de validité par défaut est de 720~minutes.

Deux points de configuration méritent attention avant une mise en production : la clé de
signature porte une valeur par défaut explicitement nommée
\code{dev-secret-key-change-me}, et le compte superutilisateur initial a un mot de passe
par défaut. Les deux doivent être redéfinis par variable d'environnement.

\subsection{Autorisation}

L'autorisation repose sur deux dépendances FastAPI centrales.
\code{require\_permission} vérifie qu'un utilisateur détient une permission et lève 403
sinon. \code{get\_scoped\_or\_404} charge une entité en appliquant simultanément le
filtre de tenant et le filtre de périmètre.

Le choix des codes de retour est délibéré. Une lecture hors périmètre répond~404, non
403 : répondre 403 confirmerait l'existence de la ressource, ce qui constitue déjà une
fuite d'information. Une écriture hors périmètre répond 403, car l'utilisateur voit
l'objet mais n'a pas le droit d'agir dessus.

\begin{remarque}[Le frontend ne protège rien]
Le masquage d'un bouton dans l'interface n'est pas un contrôle d'accès. Il améliore
l'ergonomie en n'affichant que les actions possibles. La protection réelle est
exclusivement côté serveur. Les tests d'autorisation appellent les endpoints
directement, sans passer par l'interface, précisément pour vérifier ce point.
\end{remarque}

\subsection{Isolation entre organisations}

L'identifiant d'organisation est lu dans le jeton, jamais dans le corps ou les
paramètres de la requête. Un utilisateur ne peut donc pas se déclarer membre d'une autre
organisation, même en forgeant sa requête, sans forger le jeton, ce qui exigerait la clé
de signature.

La permission \code{platform.admin} est la seule à franchir cette frontière. Elle ne
peut pas être placée par un administrateur d'entreprise dans un rôle qu'il crée : une
vérification explicite l'en empêche.

\subsection{Traçabilité}

Toute écriture significative produit une entrée dans \code{audit\_log}, conservant
l'auteur, son adresse électronique au moment de l'action, l'horodatage, le type et
l'identifiant de l'entité, un libellé lisible, l'ancienne et la nouvelle valeur en JSON,
et l'adresse IP. L'API n'expose aucune mise à jour ni suppression sur cette table : le
journal est en insertion seule. L'adresse électronique est dupliquée dans l'entrée
précisément pour que la trace reste lisible si le compte est supprimé par la suite.

% =====================================================================
%  Chapitre 4 — Interfaces
% =====================================================================
\chapter{Présentation des interfaces}
\label{chap:interfaces}

Ce chapitre parcourt les quatorze écrans de l'application dans l'ordre où une
démonstration les enchaîne naturellement. Chaque écran est décrit puis accompagné d'un
encadré \emph{Présentation orale suggérée}, directement réutilisable en soutenance.

\begin{remarque}[Insertion des captures]
Les captures sont appelées depuis le répertoire \code{rapport/images/}. Tant qu'un
fichier n'y figure pas, le document compile et affiche un cadre nommant l'image
attendue. Il suffit de déposer le fichier au nom indiqué pour que la capture apparaisse.
\end{remarque}

% =====================================================================
\section{Connexion}

L'écran de connexion est construit en deux panneaux. À gauche, sur fond vert sombre, la
promesse du produit et les quatre temps de la démarche, numérotés de
\emph{01~—~Diagnostiquer} à \emph{04~—~Piloter la valeur}. À droite, sur fond clair, le
formulaire : adresse électronique professionnelle, mot de passe avec bascule
d'affichage, et bouton de connexion doté d'un état de chargement.

Quatre états d'erreur distincts sont rendus en clair plutôt qu'en code HTTP :
identifiants invalides, compte désactivé, saisie incomplète et serveur injoignable.

Aucun identifiant de démonstration n'est écrit dans le composant React. Lorsque le
serveur tourne en mode développement, un panneau discret propose la liste des comptes,
obtenue d'un endpoint qui répond 404 en production. En production, ce panneau n'existe
donc tout simplement pas.

\capture{01-connexion.png}{Écran de connexion}{0.92}

\begin{oral}
Voici la page d'entrée de la plateforme. Le panneau de gauche résume la démarche en
quatre temps, du diagnostic au pilotage de la valeur. Le formulaire est volontairement
sobre. Un point de sécurité mérite d'être signalé : aucun compte de démonstration n'est
écrit en dur dans le code de l'interface. La liste que vous voyez ici vient d'un
endpoint qui n'existe qu'en mode développement et qui répond « non trouvé » en
production. Je me connecte avec le compte de la Direction Générale.
\end{oral}

% =====================================================================
\section{Tableau de bord Direction}

C'est l'écran d'accueil des utilisateurs porteurs de la permission
\code{dashboard.direction}. Il est composé de quatre indicateurs de tête — projets
actifs, investissement, ROI du portefeuille et risques critiques — chacun assorti d'une
précision de contexte calculée, par exemple « engagé sur 3,1~M\texteuro{} budgétés ».

Au-dessous, la matrice Impact~$\times$~Faisabilité positionne chaque opportunité, la
taille de la bulle représentant le bénéfice annuel attendu. À côté, la répartition du
portefeuille par phase TOGAF, de la qualification à la production. Enfin, un tableau des
premières priorités, cliquable jusqu'à la fiche projet.

\capture{02-dashboard-direction.png}{Tableau de bord Direction}{0.95}

\begin{oral}
Voici la vue de la Direction Générale. Ces quatre chiffres ne sont pas saisis : ils sont
agrégés à la demande depuis les projets, les coûts et les bénéfices. Le ROI de
portefeuille est le rapport entre les gains cumulés sur trois ans et le coût total de
possession. La matrice croise l'impact métier et la faisabilité ; la taille de chaque
bulle est le bénéfice annuel attendu. Le quadrant en haut à droite est ce que
l'organisation devrait lancer en premier. Toutes les lignes du tableau sont cliquables
et mènent à la fiche projet correspondante.
\end{oral}

% =====================================================================
\section{Tableau de bord Département}

Le même emplacement sert une vue différente à qui porte
\code{dashboard.department} : opportunités du périmètre, projets actifs, gains annuels
attendus et risques ouverts. Le périmètre est celui des départements de rattachement de
l'utilisateur, jamais l'organisation entière.

Suivent les opportunités du périmètre classées par score, la répartition par phase, les
projets en cours avec leur avancement et leur prochaine gate, les risques classés par
criticité et les indicateurs métier avec la part de l'écart déjà captée entre la
référence et la cible.

\capture{03-dashboard-departement.png}{Tableau de bord Département}{0.95}

\begin{oral}
Même URL, écran différent. Je suis maintenant connecté en responsable de département.
La plateforme ne se contente pas de masquer des informations : le backend filtre les
données sur les départements de rattachement de l'utilisateur. Si j'appelais l'API
directement pour lire une opportunité d'un autre département, je recevrais un 404, pas
un 403 — répondre 403 confirmerait que la ressource existe.
\end{oral}

% =====================================================================
\section{Diagnostic AS-IS et readiness}

L'écran affiche la campagne en cours, son implantation et son taux de complétion. Une
jauge donne le readiness global sur 100 et le niveau de maturité associé, de
« Inexistant » à « Optimisé ». À droite, l'échelle de maturité détaillée, puis deux
bandeaux : les dimensions bloquantes, en ambre, et les points de vigilance.

Deux colonnes séparent le readiness technique du readiness métier et organisationnel,
chacune assortie de sa moyenne sur~5. La dernière carte, « Preuves collectées »,
présente chaque dimension avec sa note, son poids, sa confiance et la preuve qui
l'étaye.

Deux actions sont proposées : exporter la campagne au format CSV et poursuivre le
diagnostic. Le second bouton s'adapte au droit de l'utilisateur : il devient « Voir le
questionnaire » pour qui n'a pas \code{diagnostic.answer}.

\capture{04-diagnostic.png}{Diagnostic AS-IS et calcul du readiness}{0.95}

\begin{oral}
Le readiness est le point de départ de la démarche. Chaque dimension est notée de 0 à 5,
pondérée par son importance et par la confiance que le répondant accorde à sa propre
réponse. Cette confiance intervient au numérateur et au dénominateur : une réponse mal
étayée ne fausse le score ni vers le haut ni vers le bas. Les deux dimensions en ambre
sont bloquantes : elles sont sous le seuil et marquées critiques, ce qui interdit un
passage en production tant qu'elles ne sont pas remédiées. Le tableau du bas montre que
chaque note est adossée à une preuve — c'est ce qui rend le diagnostic auditable.
\end{oral}

% =====================================================================
\section{Panneau de saisie du diagnostic}

Le bouton « Poursuivre le diagnostic » ouvre un panneau latéral listant les dimensions
de la campagne. Chaque ligne permet de régler la note, la confiance et la preuve. Une
dimension jamais instruite est explicitement signalée comme telle, et non affichée à
zéro.

L'enregistrement recalcule le readiness immédiatement : le score, les moyennes d'axe, le
taux de complétion et la liste des dimensions bloquantes sont rafraîchis sans
rechargement de page.

\capture{05-diagnostic-panneau.png}{Panneau de saisie des dimensions}{0.88}

\begin{oral}
Je modifie la note de la dimension « Gouvernance data » et je l'accompagne d'une preuve.
À l'enregistrement, le score global, les deux moyennes d'axe et la liste des dimensions
bloquantes sont recalculés côté serveur, puis réaffichés. Rien n'est stocké côté
navigateur : ce que vous voyez sort de la base.
\end{oral}

% =====================================================================
\section{Liste des opportunités}

L'écran recense les besoins métier. Un bandeau rappelle à l'utilisateur à périmètre
restreint qu'il ne voit que ses départements. Un second bandeau, en ambre, signale les
doublons probables détectés par le croisement interservices, en nommant les codes
concernés et les termes qu'ils partagent.

Le tableau se filtre par département et par recherche libre sur le titre ou le code.
Chaque ligne porte son type de solution recommandé, son score et ses actions. Le bouton
« Nouvelle opportunité » n'apparaît que si l'utilisateur détient
\code{opportunity.create}.

\capture{06-opportunites.png}{Liste des opportunités et détection de doublons}{0.95}

\begin{oral}
C'est ici que remontent les besoins. Le bandeau ambre est un des apports du croisement
interservices : deux départements ont exprimé un besoin très proche sans le savoir, et
la plateforme les rapproche en s'appuyant sur les termes partagés. Sans cet écran, les
deux projets seraient financés séparément.
\end{oral}

% =====================================================================
\section{Formulaire d'opportunité}

Le formulaire couvre trois blocs. Le rattachement d'abord — département et site — car le
département détermine le volet spécifique du questionnaire. Le contexte ensuite :
situation actuelle et points de douleur. Le questionnaire de qualification enfin, dont
les réponses alimentent directement les quinze faits du moteur.

À l'enregistrement, la qualification est déclenchée automatiquement. Il n'existe pas
d'opportunité non qualifiée dans la base.

\capture{07-opportunite-formulaire.png}{Création d'une opportunité}{0.92}

\begin{oral}
Le questionnaire n'est pas décoratif : chaque réponse est un fait d'entrée du moteur de
qualification. Je déclare ici des données disponibles mais non structurées, du texte
libre à traiter, et pas de génération de contenu. Je valide. Le moteur s'exécute
immédiatement et propose un type de solution. Notez qu'il n'y a pas d'étape « lancer la
qualification » : la création et la qualification sont une seule opération.
\end{oral}

% =====================================================================
\section{Architecture et données}

Cet écran tient le référentiel AS-IS en trois sections : les processus, les applications
du système d'information et les sources de données avec leur score de qualité et
l'indication d'une vérité terrain disponible.

C'est la matière première du moteur de qualification. Une source de données mal notée ou
dépourvue de vérité terrain fait échouer les règles bloquantes des types qui exigent un
apprentissage, et oriente mécaniquement la recommandation vers l'automatisation
déterministe ou la refonte de processus.

\capture{08-architecture.png}{Référentiel des processus, applications et données}{0.95}

\begin{oral}
Cet écran explique pourquoi la plateforme refuse parfois de recommander de l'IA. Si la
source de données d'un processus est notée 2 sur 5 et n'a pas de vérité terrain, la
règle bloquante du type « apprentissage supervisé » n'est pas satisfaite, et le moteur
bascule sur de l'automatisation ou sur une refonte de processus. La qualification
technique arrive donc avant la décision, pas après.
\end{oral}

% =====================================================================
\section{Portefeuille}

Le portefeuille présente toutes les opportunités classées par score, avec le bénéfice
attendu cumulé affiché en sous-titre. Cinq filtres sont disponibles : département, type
de solution, phase, score minimal et statut. Le tri est ajustable.

L'action « Recalculer la priorisation » n'est proposée qu'aux détenteurs de
\code{portfolio.prioritize}.

\capture{09-portefeuille.png}{Table du portefeuille priorisé}{0.95}

\begin{oral}
Le portefeuille est l'écran d'arbitrage. Chaque score est calculé sur six critères
pondérés, moins une pénalité de risque proportionnelle au produit probabilité par
impact. Le rang et la vague en découlent : la vague~1 est ce qu'on lance, la vague~2 ce
qu'on prépare. Rien n'est saisi à la main dans cette colonne.
\end{oral}

% =====================================================================
\section{Panneau de repriorisation}

Le panneau expose les six pondérations sous forme de curseurs. Toute grille saisie est
renormalisée à~1 avant application, de sorte qu'un utilisateur qui pousse un seul critère
obtient une pondération valide sans avoir à ajuster les cinq autres.

Le recalcul réévalue tout le portefeuille, réattribue les rangs et les vagues, et
journalise l'opération dans le journal d'audit avec l'ancienne et la nouvelle
pondération.

\capture{10-repriorisation.png}{Ajustement des pondérations de priorisation}{0.88}

\begin{oral}
C'est l'écran qui montre que la priorisation est un outil de pilotage, pas une formule
figée. Si la Direction décide que le contexte impose de privilégier l'impact financier,
elle pousse ce curseur et relance le calcul. Le classement change réellement : sur le
jeu de démonstration, douze opportunités changent de rang. Et le changement de
pondération est lui-même tracé dans le journal d'audit, avec l'ancienne et la nouvelle
grille.
\end{oral}

% =====================================================================
\section{Liste des fiches projet}

Un écran de sélection listant les fiches disponibles, chacune renvoyant au dossier de
cadrage correspondant. Les fiches ne sont pas des documents importés : elles sont
assemblées à la demande à partir des données vivantes du portefeuille.

\capture{11-fiches.png}{Sélection d'une fiche projet}{0.9}

\begin{oral}
Ces fiches ne sont pas stockées comme des documents. Elles sont reconstruites à chaque
ouverture à partir de l'opportunité, de sa recommandation, de son score, de ses
indicateurs, de ses risques et de son business case. Une fiche ne peut donc pas être
périmée par rapport aux données.
\end{oral}

% =====================================================================
\section{Fiche projet}

C'est l'écran le plus dense de l'application, et le point d'aboutissement de la
démarche. Un bandeau d'identité — département, sponsor, responsable métier, responsable
technique, phase et statut — précède deux colonnes.

La colonne principale déroule le contexte AS-IS et ses points de douleur, la solution
recommandée avec son patron de solution, sa justification rédigée, les règles satisfaites
affichées en pastilles, la première alternative écartée et les garde-fous à mettre en
place ; puis les indicateurs et le bénéfice attendu, et enfin les risques et la
conformité.

La colonne latérale, sur fond sombre, affiche le score de priorité en grand, le rang sur
la taille du portefeuille, la vague, puis la contribution détaillée de chacun des six
critères et la pénalité de risque. Viennent ensuite le business case, réservé aux
porteurs de \code{finance.read}, la décision de comité et un bloc d'explicabilité.

\capture{12-fiche-projet.png}{Fiche projet complète}{0.95}

\begin{oral}
Cette fiche est ce qu'un comité d'investissement reçoit. Trois choses la distinguent
d'un document classique. D'abord, la recommandation est justifiée : les pastilles vertes
sont les règles qui ont été satisfaites, et la pastille grise est l'alternative écartée
avec le motif. Ensuite, le score de la colonne de droite est décomposé critère par
critère, y compris la pénalité de risque. Enfin, les garde-fous sont générés
automatiquement — sur une opportunité en périmètre GxP, la plateforme ajoute la
validation informatique, la traçabilité ALCOA+ et la revue humaine obligatoire.
\end{oral}

\begin{ecart}
Le bouton porte le libellé « Export Word / PDF ». L'API produit en réalité un rendu JSON
et une page HTML mise en page pour l'impression : l'utilisateur obtient un PDF par la
fonction d'impression de son navigateur, mais aucun fichier \code{.docx} n'est généré.
Le libellé est en avance sur la fonction.
\end{ecart}

% =====================================================================
\section{Valeur et ROI}

L'écran consolide les business cases sur un horizon de trois ans. Il présente les gains
annuels par département, le classement des projets par ROI pour les huit premiers, le
suivi des indicateurs avant et après, et un tableau détaillé par projet donnant le coût
total de possession, les gains annuels, le ROI et le délai de retour.

\capture{13-valeur.png}{Pilotage de la valeur et du retour sur investissement}{0.95}

\begin{oral}
C'est l'écran qui ferme la boucle. Le business case ne sert plus seulement à obtenir un
budget : la colonne des indicateurs compare la valeur promise à la valeur effectivement
mesurée. Le délai de retour se calcule sur le seul investissement initial rapporté au
gain annuel net d'exploitation — y faire figurer trois années de coûts récurrents qui ne
sont pas encore engagés allongerait artificiellement le délai.
\end{oral}

% =====================================================================
\section{Administration}

L'écran d'administration est organisé en trois onglets, dont l'affichage dépend des
droits : \emph{Organisation}, \emph{Utilisateurs} et \emph{Rôles \& permissions}.

L'onglet Organisation regroupe les informations de l'entreprise, ses sites et ses
départements. L'onglet Utilisateurs permet de créer un compte, de l'activer ou de le
désactiver, de lui attribuer des rôles et de le rattacher à des départements et à des
sites. L'onglet Rôles affiche la grille rôle~$\times$~permission ; les rôles fournis par
la plateforme y sont en lecture seule, un rôle personnalisé devant être créé pour
ajuster les droits.

\capture{14-administration.png}{Console d'administration, onglet Rôles et permissions}{0.95}

\begin{oral}
Voici le module M1. Les onze rôles que vous voyez ne sont pas des libellés décoratifs :
chacun est un jeu de permissions parmi les quarante et une du catalogue, associé à un
niveau de périmètre. Les rôles système ne sont pas modifiables, pour qu'un
administrateur ne puisse pas se retirer involontairement des droits dont la plateforme
dépend. Il peut en revanche créer ses propres rôles. Une seule permission lui est
interdite : celle qui franchit la frontière entre organisations.
\end{oral}

% =====================================================================
\section{Journal d'audit}

Le dernier écran présente les événements en insertion seule. Chaque ligne porte
l'horodatage, l'auteur et son adresse électronique, l'action, l'entité concernée et son
libellé. Le détail expose l'ancienne et la nouvelle valeur.

Le sous-titre de l'écran énonce la propriété qui fait sa valeur : aucune entrée ne peut
être modifiée ni supprimée. L'API n'expose ni mise à jour ni suppression sur cette
table.

\capture{15-audit.png}{Journal d'audit en insertion seule}{0.95}

\begin{oral}
C'est l'écran qui répond à la question « pourquoi cette décision ? » six mois plus tard.
Chaque modification significative laisse une trace avec son auteur, sa date, l'ancienne
et la nouvelle valeur. Le point important est qu'il n'existe aucun moyen, ni par
l'interface ni par l'API, de modifier ou de supprimer une entrée : le journal est en
insertion seule. C'est une exigence directe du contexte réglementaire pharmaceutique.
\end{oral}

% =====================================================================
%  Chapitre 5 — Tests et validation
% =====================================================================
\chapter{Tests et validation}
\label{chap:tests}

\section{Stratégie de validation}

La validation s'appuie sur trois niveaux, choisis en fonction de ce qu'il s'agit de
prouver.

\textbf{Les services métier se testent sans serveur.} Les formules de readiness, de
qualification, de priorisation et de valorisation sont des fonctions sans effet de bord.
Elles se vérifient par comparaison de valeurs exactes, sans base de données ni requête
HTTP.

\textbf{Les règles d'autorisation se testent par l'API.} Un contrôle d'accès ne se
prouve pas en lisant le code de l'interface. Les tests d'autorisation construisent un
jeton pour un rôle donné, appellent l'endpoint directement et vérifient le code de
retour. C'est la seule manière de garantir qu'un utilisateur qui contourne l'interface
se heurte bien au serveur.

\textbf{Les parcours d'écran se vérifient dans un navigateur réel.} Les enchaînements
d'interface — connexion selon le rôle, ouverture des panneaux latéraux, recalcul de la
priorisation, export CSV — ont été exécutés dans un navigateur automatisé, avec capture
à chaque étape.

\section{Composition de la suite automatisée}

La suite compte \textbf{112~tests} répartis en sept fichiers.

\begin{table}[H]
\centering
\caption{Répartition de la suite de tests}
\label{tab:repartition-tests}
\small
\begin{tabular}{@{}L{5.2cm} C{1.6cm} L{7.2cm}@{}}
\toprule
\textbf{Fichier} & \textbf{Tests} & \textbf{Objet} \\
\midrule
\code{test\_permissions.py}        & 40 & Comportement de chacun des 11 rôles \\
\code{test\_login.py}              & 19 & Authentification et aide de développement \\
\code{test\_opportunities.py}      & 15 & Cycle de vie, qualification, priorisation \\
\code{test\_tenant\_isolation.py}  & 13 & Étanchéité entre organisations \\
\code{test\_diagnostic.py}         & 11 & Sémantique du readiness \\
\code{test\_scope.py}              & 10 & Bornage par département \\
\code{test\_endpoint\_protection.py}&  4 & Garde-fous structurels sur l'API \\
\midrule
\textbf{Total}                     & \textbf{112} & \\
\bottomrule
\end{tabular}
\end{table}

\section{Résultats}

\begin{table}[H]
\centering
\caption{Résultat de l'exécution de la suite}
\label{tab:resultats-tests}
\small
\begin{tabular}{@{}L{5.6cm} L{5.4cm} C{2.8cm}@{}}
\toprule
\textbf{Vérification} & \textbf{Objectif} & \textbf{Résultat} \\
\midrule
Suite backend complète & 112 tests & \fait \ 112 réussis \\
Durée d'exécution & Boucle de retour courte & 12,15 s \\
Typage du frontend & \code{tsc --noEmit} sans erreur & \fait \ aucune erreur \\
\bottomrule
\end{tabular}
\end{table}

\section{Détail des vérifications}

\subsection{Sémantique du readiness}

\begin{longtable}{@{}L{7.4cm} L{5.2cm} C{1.5cm}@{}}
\caption{Tests du calcul de readiness}
\label{tab:tests-readiness}\\
\toprule
\textbf{Test} & \textbf{Ce qu'il garantit} & \textbf{Résultat} \\
\midrule
\endfirsthead
\toprule
\textbf{Test} & \textbf{Ce qu'il garantit} & \textbf{Résultat} \\
\midrule
\endhead
\bottomrule
\endfoot
Le taux de complétion est calculé, non saisi
& Aucune valeur de complétion ne peut être écrite depuis l'API & OK \\
Une dimension non instruite ne compte pas comme un zéro
& Le score n'est pas mécaniquement dégradé par une absence de réponse & OK \\
Les moyennes d'axe ignorent les dimensions non instruites
& Les deux moyennes restent comparables entre campagnes & OK \\
Instruire une dimension fait progresser la complétion
& Le lien réponse $\rightarrow$ complétion est effectif & OK \\
Instruire une dimension modifie le score
& Le score est bien recalculé, non mis en cache & OK \\
Une dimension critique sous le seuil devient bloquante
& La règle bloquante / vigilance est appliquée & OK \\
Répondre au diagnostic exige la permission
& L'écriture est refusée sans \code{diagnostic.answer} & OK \\
La campagne d'un autre tenant est introuvable
& Réponse 404, pas 403 & OK \\
\end{longtable}

\subsection{Autorisation par rôle}

Quarante tests vérifient que chaque rôle fait ce qu'il doit faire \emph{et} échoue sur
ce qu'il ne doit pas faire. La seconde moitié est la plus importante : elle constitue la
preuve que les permissions sont réellement discriminantes.

\begin{longtable}{@{}L{7.4cm} L{5.2cm} C{1.5cm}@{}}
\caption{Extraits des tests d'autorisation}
\label{tab:tests-permissions}\\
\toprule
\textbf{Test} & \textbf{Ce qu'il garantit} & \textbf{Résultat} \\
\midrule
\endfirsthead
\toprule
\textbf{Test} & \textbf{Ce qu'il garantit} & \textbf{Résultat} \\
\midrule
\endhead
\bottomrule
\endfoot
Le catalogue ne contient aucun code inconnu
& Aucune permission orpheline dans la base & OK \\
Chaque modèle de rôle ne référence que des permissions existantes
& Pas de rôle pointant vers une permission supprimée & OK \\
Un administrateur d'entreprise ne peut pas s'octroyer le rôle de plateforme
& La frontière entre tenants n'est pas franchissable & OK \\
\ldots{} ni créer un rôle portant une permission de plateforme
& L'élévation de privilège par rôle personnalisé est bloquée & OK \\
Un rôle système n'est pas modifiable par un administrateur d'entreprise
& Les rôles livrés restent cohérents & OK \\
La Direction ne crée pas d'utilisateur
& Consulter et arbitrer n'est pas administrer & OK \\
Finance ne modifie pas la recommandation, ne pilote pas les projets
& Séparation des responsabilités & OK \\
L'AI Lead requalifie une opportunité, n'administre pas les comptes
& \code{recommendation.override} est exclusive à ce rôle & OK \\
Le Data Office ne gère pas le référentiel applicatif
& Data et architecture sont deux périmètres distincts & OK \\
Le PMO ne valide pas les bénéfices financiers
& La validation financière reste chez Finance et la Direction & OK \\
L'auditeur ne peut rien écrire, modifier ni supprimer
& Lecture seule effective sur tous les modules & OK \\
Le contributeur ne voit ni le portefeuille consolidé ni les données financières
& Le profil le plus restreint l'est réellement & OK \\
\end{longtable}

\subsection{Périmètre et étanchéité}

Dix tests couvrent le bornage par département et treize l'étanchéité entre
organisations. Ces derniers construisent une seconde organisation complète, avec ses
propres utilisateurs et rôles, et vérifient qu'aucune lecture, écriture, suppression,
affectation de département, attribution de rôle ni entrée d'audit ne traverse la
frontière. Le test final vérifie que toutes les listes du second locataire sont vides
des données du premier.

Un point mérite d'être explicité : un utilisateur d'une autre organisation qui appelle
une ressource existante reçoit \textbf{404}, non 403. Répondre 403 confirmerait
l'existence de la ressource, ce qui est déjà une fuite d'information. En revanche une
écriture hors périmètre \emph{à l'intérieur} de sa propre organisation reçoit bien 403,
puisque l'utilisateur voit l'objet mais n'a pas le droit d'agir dessus.

\subsection{Garde-fous structurels}

Quatre tests ne vérifient pas un comportement mais une propriété du code lui-même. Ils
parcourent la table de routage de l'application et contrôlent que :

\begin{itemize}
  \item aucun endpoint n'est exposé sans contrôle, à l'exception d'une liste explicite
    de routes publiques, chacune assortie d'un commentaire justifiant son ouverture ;
  \item aucune dépendance d'autorisation n'a été silencieusement transformée en
    paramètre de requête ou en corps de requête ;
  \item toutes les routes métier vivent sous le préfixe versionné.
\end{itemize}

Le deuxième de ces tests a été ajouté après un incident réel. Un import manquant avait
transformé une annotation de périmètre en paramètre de requête : FastAPI répondait 422
avec un message obscur, et surtout le contrôle de périmètre n'était plus appliqué. Le
test rend cette classe d'erreur impossible à réintroduire sans être détectée.

\section{Vérification des interfaces}

Les parcours suivants ont été exécutés dans un navigateur automatisé, avec capture à
chaque étape.

\begin{table}[H]
\centering
\caption{Parcours d'interface vérifiés}
\label{tab:tests-ui}
\small
\begin{tabular}{@{}L{5.6cm} L{6.4cm} C{1.8cm}@{}}
\toprule
\textbf{Parcours} & \textbf{Attendu} & \textbf{Résultat} \\
\midrule
Connexion par rôle
& Chaque rôle atterrit sur son écran d'accueil pertinent & OK \\
Navigation latérale
& Le menu ne contient que les écrans accessibles & OK \\
Panneau de diagnostic
& La saisie recalcule score, moyennes et complétion & OK \\
Recalcul de la priorisation
& Les curseurs modifient réellement le classement & OK \\
Création d'opportunité
& La qualification est déclenchée à l'enregistrement & OK \\
Export CSV du diagnostic
& Fichier téléchargé, accents corrects, résumé inclus & OK \\
Affichage mobile
& Aucun débordement horizontal jusqu'à 390~px & OK \\
\bottomrule
\end{tabular}
\end{table}

\section{Jeu de démonstration}

Les tests s'appuient sur des données créées à la volée ; la démonstration s'appuie sur
un jeu ensemencé, calibré pour être crédible et cohérent avec les maquettes.

\begin{table}[H]
\centering
\caption{Contenu du jeu de démonstration}
\label{tab:seed}
\small
\begin{tabular}{@{}L{6cm} C{2.4cm} L{5.4cm}@{}}
\toprule
\textbf{Élément} & \textbf{Volume} & \textbf{Remarque} \\
\midrule
Organisation & 1 & Laboratoire pharmaceutique \\
Sites & 4 & Tunis, Lyon, Casablanca, Paris \\
Départements & 8 & Du Commercial aux Affaires réglementaires \\
Utilisateurs & 13 & Au moins un par rôle fonctionnel \\
Réponses de diagnostic & 16 & Deux axes, avec preuves et confiances \\
Opportunités & 18 & Toutes qualifiées et scorées \\
Projets & 13 & Répartis sur les cinq phases \\
Entrées d'audit & 440 au seed & Volume qui croît ensuite à chaque écriture \\
\bottomrule
\end{tabular}
\end{table}

Les valeurs calculées à partir de ce jeu sont stables et reproductibles : readiness
global de 62,0 sur 100, moyennes d'axe de 2,81 et 3,12 sur~5, deux dimensions bloquantes
— gouvernance des données et infrastructure IA —, quatre premières priorités à 87, 81,
74 et 69~points, ROI de portefeuille de 183,9~\% et délai de retour moyen de 11~mois.

Ces valeurs ne sont écrites nulle part : elles résultent des formules appliquées aux
données ensemencées. C'est la raison pour laquelle le jeu a été calibré avec soin — une
démonstration dont les chiffres ne tiennent pas ensemble se voit immédiatement.

\begin{remarque}[Sur les comptes de démonstration]
Le mot de passe commun aux comptes du seed n'est pas écrit dans le code de l'interface.
Il est lu d'une variable d'environnement et n'est publié que par un endpoint d'aide qui
répond 404 dès que l'environnement n'est pas déclaré comme développement. Ces comptes
restent strictement destinés au développement et à la démonstration.
\end{remarque}

\section{Difficultés rencontrées}

Les difficultés ci-dessous ont été effectivement rencontrées pendant le développement et
sont documentées ici parce qu'elles ont conduit à des corrections vérifiables dans le
code.

\textbf{Une dimension ajoutée au référentiel dégradait tous les scores.} Les dimensions
non encore instruites étaient comptées comme des zéros. Ajouter une question au
questionnaire faisait donc chuter le readiness de toutes les campagnes existantes,
alors que rien n'avait changé sur le terrain. La notion de dimension instruite, adossée
à une confiance strictement positive, a résolu le problème : une dimension non mesurée
sort du calcul et pèse seulement sur le taux de complétion.

\textbf{Une gate franchie faisait reculer un projet.} Valider une gate déjà franchie
ramenait le projet à une phase antérieure. Un garde-fou répond désormais 409 avec un
message explicite selon que la gate est déjà franchie ou pas encore atteinte.

\textbf{Le recalcul des synergies écrasait les valeurs ensemencées.} La priorisation
recalculait automatiquement les synergies, ce qui modifiait les scores à chaque passage
et rendait la démonstration instable. Le recalcul est devenu explicite plutôt
qu'automatique.

\textbf{Le délai de retour ne concordait pas avec les hypothèses.} La formule initiale
rapportait le coût total de possession au gain mensuel, ce qui faisait porter au délai
de retour trois années de coûts récurrents non encore engagés. Elle a été corrigée pour
rapporter le seul investissement initial au gain annuel net d'exploitation.

\textbf{Un contrôle de périmètre avait disparu silencieusement.} Un import manquant
transformait une annotation en paramètre de requête ; FastAPI répondait 422 et le
contrôle n'était plus exécuté. Au-delà de la correction, un test structurel garantit
désormais qu'aucune dépendance ne peut subir ce sort sans être détectée.

\textbf{Un test passait seul et échouait en suite.} Un test d'interface laissait des
pondérations modifiées en base, ce qui faussait le test de priorisation exécuté ensuite.
Le montage de test rétablit désormais les pondérations par défaut avant et après.

\textbf{Le développement sous Windows a posé ses propres questions.} La documentation
d'installation décrivait l'activation d'environnement virtuel en syntaxe POSIX, ce qui
produisait une erreur sous PowerShell. Les deux fichiers de documentation donnent
désormais la commande adaptée à chaque terminal, avec une note sur la politique
d'exécution.

% =====================================================================
%  Conclusion générale
% =====================================================================
\chapter*{Conclusion générale}
\addcontentsline{toc}{chapter}{Conclusion générale}
\markboth{Conclusion générale}{}
\label{chap:conclusion}

\section*{Bilan du travail réalisé}

Le projet a abouti à une plateforme fonctionnelle couvrant le parcours complet de
l'adoption de l'intelligence artificielle dans une organisation pharmaceutique, du
diagnostic de maturité à la mesure de la valeur captée.

En chiffres, la réalisation représente 34~tables, 99~endpoints répartis en dix domaines
fonctionnels, sept services métier isolés du transport HTTP, trois migrations de schéma,
quatorze écrans, 41~permissions, onze rôles, cinq niveaux de périmètre et 112~tests
automatisés qui passent.

Neuf des douze modules du cahier des charges sont pleinement implémentés. Trois le sont
partiellement, et le chapitre~\ref{chap:analyse} en détaille l'écart sans le minimiser :
l'export bureautique de la fiche projet produit du HTML imprimable plutôt qu'un fichier
\code{.docx}, le suivi de projet dispose de son API mais pas de son écran dédié, et les
notifications sont persistées sans panneau de consultation.

Trois choix de conception structurent l'ensemble.

\textbf{L'explicabilité comme contrainte de départ, non comme option.} Chaque calcul
renvoie sa décomposition : le readiness rend ses dimensions avec leur preuve, le score
de priorité rend la contribution de ses six critères, la recommandation rend ses règles
satisfaites, ses alternatives écartées et ses garde-fous. Cette contrainte a orienté le
modèle de données — les explications sont persistées en JSON — autant que les formules.

\textbf{La sécurité côté serveur, exclusivement.} L'interface adapte ce qu'elle affiche
aux droits de l'utilisateur, mais aucune décision d'accès ne lui est déléguée. Les
tests d'autorisation appellent l'API directement, sans passer par l'interface, parce que
c'est ainsi qu'un attaquant procéderait.

\textbf{La séparation entre calcul et exposition.} Aucune formule ne vit dans un router.
Cette discipline a rendu les algorithmes testables sans serveur ni base de données, et
c'est ce qui a permis de les corriger tard dans le projet sans crainte de régression.

\section*{Limites de la version actuelle}

Les limites ci-dessous sont réelles et constatées dans le code livré.

\textbf{Aucun modèle d'apprentissage.} La plateforme recommande des solutions
d'intelligence artificielle sans en contenir aucune. Le moteur de qualification est un
système de règles pondérées : il ne découvre rien qu'un expert n'ait prévu et n'apprend
pas des décisions passées. Ce choix est argumenté à la section~\ref{sec:pas-dia} ; il
n'en reste pas moins une limite.

\textbf{La détection de doublons est lexicale.} Le croisement interservices repère les
termes partagés entre deux expressions de besoin. Deux besoins identiques formulés avec
des vocabulaires métier différents ne sont pas rapprochés.

\textbf{Le suivi de projet n'a pas d'écran.} Phases, jalons, tâches, risques,
mitigations, incidents et passage de gates sont pilotables par l'API, pas par
l'interface. Le PMO retrouve l'avancement et la prochaine gate dans son tableau de bord,
ce qui couvre le pilotage mais pas la gestion détaillée.

\textbf{L'export bureautique n'existe pas.} Le libellé du bouton est en avance sur la
fonction.

\textbf{Pas de réinitialisation de mot de passe autonome.} Un utilisateur qui oublie son
mot de passe doit s'adresser à l'administrateur de son organisation.

\textbf{Pas de rafraîchissement de jeton.} Le jeton est valable douze heures, après quoi
l'utilisateur se reconnecte. Il n'existe ni jeton de rafraîchissement, ni liste de
révocation : un jeton volé reste exploitable jusqu'à son expiration. La seule parade
actuelle est la désactivation du compte, qui est bien effective à chaque requête — le
serveur vérifie l'état du compte après avoir validé le jeton — mais qui suppose de
désactiver l'utilisateur légitime.

\textbf{Configuration par défaut à durcir.} La clé de signature et le mot de passe du
superutilisateur initial portent des valeurs par défaut qui doivent impérativement être
redéfinies avant toute mise en production.

\textbf{Écart d'environnement.} Le développement s'est fait sur PostgreSQL~12.4 en
local, tandis que le fichier Docker déclare PostgreSQL~16. Aucune fonctionnalité
spécifique à la version~16 n'est utilisée, mais l'écart n'a pas été validé en
conditions réelles.

\section*{Perspectives}

Les évolutions ci-dessous sont des propositions. Aucune n'est implémentée.

\textbf{Assistance par modèle de langage, en entrée du moteur.} Un modèle pourrait
pré-remplir les quinze faits de qualification à partir d'une description libre du
besoin, rédiger un premier jet de fiche projet et détecter des doublons sémantiques
là où la comparaison lexicale échoue. Le principe directeur serait de laisser le modèle
\emph{proposer des entrées} et le moteur de règles \emph{décider} : la traçabilité de la
décision resterait intacte, et le gain se situerait sur la saisie, qui est aujourd'hui
le point de friction.

\textbf{Apprentissage sur l'historique des décisions.} Après plusieurs dizaines de
cycles complets — qualification, décision de comité, résultat mesuré — l'organisation
disposerait enfin d'un corpus. Les pondérations de priorisation pourraient alors être
calibrées sur les projets qui ont réellement tenu leurs promesses, plutôt que fixées
a~priori.

\textbf{Écran de pilotage de projet.} Exposer l'API de suivi déjà en place :
diagramme de Gantt des jalons, tableau des tâches, matrice des risques et parcours de
gates.

\textbf{Export bureautique réel.} Générer un fichier \code{.docx} ou un PDF mis en page
côté serveur, pour aligner la fonction sur le libellé.

\textbf{Notifications actionnables.} Panneau de consultation, marquage lu et non lu, et
déclenchement sur les événements qui appellent une action : gate à franchir, indicateur
qui décroche, décision de comité en attente.

\textbf{Durcissement de l'authentification.} Jeton de rafraîchissement, révocation côté
serveur, réinitialisation autonome du mot de passe et, en périmètre réglementé,
authentification à deux facteurs.

\textbf{Exploitation des périmètres par site.} Le modèle porte déjà le rattachement des
utilisateurs aux sites et le niveau de périmètre correspondant. Les filtres appliqués
aujourd'hui s'arrêtent au département ; étendre le bornage au site serait une extension
naturelle du dispositif existant.

\section*{Apport personnel}

Ce projet a demandé de tenir ensemble trois exigences qui se contredisent
régulièrement : la richesse fonctionnelle, la rigueur du contrôle d'accès et
l'honnêteté sur ce que le système fait réellement.

La dernière est celle qui a le plus appris. Il aurait été facile de présenter le moteur
de qualification comme une intelligence artificielle : le vocabulaire du domaine s'y
prête, et personne n'aurait vérifié. Choisir de le nommer pour ce qu'il est — un système
expert à base de règles — et d'argumenter pourquoi ce choix est adapté au problème posé
s'est révélé plus solide que l'inverse. Un moteur explicable se défend devant un comité,
devant un inspecteur et devant un jury ; un modèle qu'on ne sait pas justifier ne se
défend nulle part.


% =====================================================================
%  Annexes
% =====================================================================
\appendix
% =====================================================================
%  Annexe A — Guide de soutenance
% =====================================================================
\chapter{Guide de soutenance}
\label{ann:soutenance}

Cette annexe rassemble ce qui est nécessaire le jour de la présentation : les comptes à
utiliser, un déroulé de démonstration minuté, et un jeu de questions probables avec des
réponses appuyées sur le code réel.

% =====================================================================
\section{Préparation}

\subsection{Lancement de l'environnement}

Trois terminaux, dans cet ordre.

\begin{enumerate}
  \item \textbf{Base de données.} Vérifier que PostgreSQL est démarré et que la base
    \code{medicacom} existe.
  \item \textbf{Backend.} Depuis \code{backend/}, activer l'environnement virtuel puis
    lancer le serveur sur le port~8000. Sous PowerShell, l'activation se fait avec
    \code{.\textbackslash.venv\textbackslash Scripts\textbackslash Activate.ps1} ; sous
    Git~Bash, avec \code{source .venv/Scripts/activate}. La confusion entre les deux
    syntaxes est la cause la plus fréquente de l'erreur
    « \code{uvicorn} n'est pas reconnu ».
  \item \textbf{Frontend.} Depuis \code{frontend/}, \code{npm run dev}. L'interface
    répond sur le port~5173.
\end{enumerate}

Ouvrir un onglet supplémentaire sur \code{http://localhost:8000/docs}. La documentation
interactive de l'API est utile si le jury demande une preuve côté serveur.

\subsection{Comptes de démonstration}

Tous les comptes ci-dessous partagent le mot de passe défini par la variable
\code{DEMO\_PASSWORD}. Ils sont réservés au développement et à la démonstration.

\begin{table}[H]
\centering
\caption{Comptes du jeu de démonstration}
\label{tab:comptes}
\small
\begin{tabular}{@{}L{5.3cm} L{3.4cm} L{5.1cm}@{}}
\toprule
\textbf{Adresse} & \textbf{Rôle} & \textbf{Écran d'accueil} \\
\midrule
admin@medicacom.com            & Admin. Medicacom & Administration \\
admin.entreprise@medicacom.com & Admin. entreprise & Administration \\
marwen.b@medicacom.com         & Direction Générale & Tableau de bord Direction \\
l.trabelsi@medicacom.com       & Resp. département & Tableau de bord Département \\
h.gharbi@medicacom.com         & Resp. département & Tableau de bord Département \\
c.lambert@medicacom.com        & Resp. département & Tableau de bord Département \\
a.rousseau@medicacom.com       & Contributeur métier & Opportunités \\
k.zouari@medicacom.com         & DSI / Architecte & Architecture \& données \\
n.hamdi@medicacom.com          & Data Office & Architecture \& données \\
y.benali@medicacom.com         & AI Lead & Portefeuille \\
m.dupuis@medicacom.com         & Finance & Valeur \& ROI \\
f.moreau@medicacom.com         & PMO & Portefeuille \\
audit@medicacom.com            & Auditeur & Portefeuille \\
\bottomrule
\end{tabular}
\end{table}

% =====================================================================
\section{Déroulé de démonstration (10 à 15 minutes)}

\subsection*{Séquence 1 — Connexion et parcours (1 min 30)}

\textbf{Où :} page de connexion.\quad\textbf{Quoi :} se connecter avec
\code{marwen.b@medicacom.com}.

\begin{oral}
Medicacom AOI Transformation répond à un problème simple à énoncer et difficile à
résoudre : dans un laboratoire pharmaceutique, les idées d'IA ne manquent pas, mais rien
ne permet de décider lesquelles financer. La plateforme couvre le parcours complet, du
diagnostic de maturité à la mesure de la valeur réellement captée. Je me connecte en
Direction Générale.
\end{oral}

\subsection*{Séquence 2 — Tableau de bord Direction (1 min 30)}

\textbf{Où :} écran d'accueil.\quad\textbf{Quoi :} montrer les quatre indicateurs, puis
la matrice Impact~$\times$~Faisabilité.

\begin{oral}
Ces quatre chiffres ne sont pas saisis. Ils sont agrégés à la demande depuis les
projets, les coûts et les bénéfices. Le ROI de portefeuille rapporte les gains cumulés
sur trois ans au coût total de possession. Dans la matrice, chaque bulle est une
opportunité ; sa taille est le bénéfice annuel attendu, et le quadrant supérieur droit
est ce que l'organisation devrait lancer en premier.
\end{oral}

\subsection*{Séquence 3 — Diagnostic et readiness (2 min)}

\textbf{Où :} menu \emph{Diagnostic}.\quad\textbf{Quoi :} montrer la jauge, les deux
bandeaux d'alerte, puis descendre jusqu'au tableau des preuves.

\begin{oral}
Le readiness est le point de départ. Chaque dimension est notée de 0 à 5, pondérée par
son importance et par la confiance que le répondant accorde à sa propre réponse. Cette
confiance intervient au numérateur et au dénominateur, donc une réponse mal étayée ne
fausse le score dans aucun sens. Les deux dimensions en ambre sont bloquantes : sous le
seuil et marquées critiques, elles interdisent un passage en production. Et chaque note
est adossée à une preuve — c'est ce qui rend le diagnostic auditable.
\end{oral}

\textbf{Action facultative :} ouvrir « Poursuivre le diagnostic », modifier une note,
enregistrer, montrer que le score bouge.

\subsection*{Séquence 4 — Opportunité et qualification (2 min 30)}

\textbf{Où :} menu \emph{Opportunités}, puis « Nouvelle opportunité ».\quad
\textbf{Quoi :} d'abord signaler le bandeau de doublons, puis créer une opportunité.

\begin{oral}
Le bandeau ambre est un apport du croisement interservices : deux départements ont
exprimé un besoin très proche sans le savoir, et la plateforme les rapproche sur leurs
termes partagés. Sans cet écran, les deux projets seraient financés séparément.

Je crée maintenant une opportunité. Ce questionnaire n'est pas décoratif : chaque
réponse est un fait d'entrée du moteur. Je déclare des données disponibles mais non
structurées, du texte libre à traiter, pas de génération de contenu. Je valide — et la
qualification s'exécute immédiatement. Il n'y a pas d'étape séparée : dans cette
plateforme, il n'existe pas d'opportunité non qualifiée.
\end{oral}

\subsection*{Séquence 5 — Portefeuille et repriorisation (2 min)}

\textbf{Où :} menu \emph{Portefeuille}, puis « Recalculer la priorisation ».

\begin{oral}
Le score de chaque ligne combine six critères pondérés, moins une pénalité de risque
proportionnelle au produit probabilité par impact. Le rang et la vague en découlent.

Si la Direction décide que le contexte impose de privilégier l'impact financier, elle
pousse ce curseur et relance le calcul. Je le fais : douze opportunités changent de
rang. Et ce changement de pondération est lui-même tracé dans le journal d'audit, avec
l'ancienne et la nouvelle grille.
\end{oral}

\subsection*{Séquence 6 — Fiche projet (2 min 30)}

\textbf{Où :} cliquer sur la première ligne du portefeuille.

\begin{oral}
C'est le livrable que reçoit un comité d'investissement. Trois choses le distinguent
d'une présentation classique. D'abord, la recommandation est justifiée : les pastilles
sont les règles satisfaites, et l'alternative écartée est nommée avec son motif.
Ensuite, le score dans la colonne de droite est décomposé critère par critère, pénalité
de risque comprise. Enfin, les garde-fous sont générés automatiquement : sur une
opportunité en périmètre GxP, la plateforme ajoute la validation informatique, la
traçabilité ALCOA+ et la revue humaine obligatoire.
\end{oral}

\subsection*{Séquence 7 — Changement de rôle (2 min)}

\textbf{Où :} se déconnecter, se reconnecter avec \code{l.trabelsi@medicacom.com}.

\begin{oral}
Même application, même URL d'accueil, écran différent. Le menu latéral a changé aussi :
il ne contient que les écrans auxquels ce rôle a droit.

Mais l'essentiel n'est pas ce qu'on voit. Masquer un bouton dans React n'est pas un
contrôle d'accès. Si j'appelais l'API directement pour lire une opportunité d'un autre
département, je recevrais 404 — et non 403, parce que répondre 403 confirmerait que la
ressource existe. Quarante tests vérifient ce comportement rôle par rôle, en appelant
l'API sans passer par l'interface.
\end{oral}

\subsection*{Séquence 8 — Valeur et audit (1 min 30)}

\textbf{Où :} menus \emph{Valeur \& ROI} puis \emph{Journal d'audit}.

\begin{oral}
L'écran Valeur ferme la boucle : le business case ne sert plus seulement à obtenir un
budget, il compare la valeur promise à la valeur mesurée.

Et le journal répond à la question « pourquoi cette décision ? » six mois plus tard.
Chaque modification laisse son auteur, sa date, l'ancienne et la nouvelle valeur. Point
important : il n'existe aucun moyen, ni par l'interface ni par l'API, de modifier ou de
supprimer une entrée. Le journal est en insertion seule, ce qui est une exigence directe
du contexte pharmaceutique.
\end{oral}

\subsection*{Conclusion orale (30 s)}

\begin{oral}
Un dernier point, qui est aussi le plus important. Le produit s'appelle « AOI
Transformation » et recommande des solutions d'intelligence artificielle. Il n'en
contient aucune. Le moteur qui produit ces recommandations est un système à base de
règles pondérées, écrites par un expert, entièrement déterministe. C'est un choix
assumé : sans historique de décisions, un modèle n'aurait rien à apprendre, et en
périmètre réglementé une décision non explicable est une décision indéfendable.
L'intégration d'un modèle de langage figure dans les perspectives du rapport, pas dans
les fonctionnalités livrées.
\end{oral}

% =====================================================================
\section{Questions probables du jury}

\subsection*{Sur l'intelligence artificielle}

\textbf{1. Quel modèle d'IA avez-vous utilisé ?}\\
Aucun. La plateforme ne contient ni modèle entraîné, ni appel à un fournisseur externe.
Ce que le nom du produit désigne, c'est le domaine traité — la gouvernance de l'adoption
de l'IA — pas la technologie employée. Le moteur qui produit les recommandations est un
système expert à base de 39~règles pondérées réparties sur 8~types de solution.

\textbf{2. Comment pouvez-vous l'affirmer ?}\\
Ni \code{requirements.txt} ni \code{package.json} ne déclarent de bibliothèque
d'apprentissage ou de client de fournisseur de modèle. Aucun fichier de modèle entraîné
n'existe dans le dépôt, aucune clé d'API n'est attendue par la configuration, et le
backend n'émet aucune requête HTTP sortante. Le seul client HTTP installé sert aux
tests.

\textbf{3. Pourquoi ne pas avoir intégré un modèle ?}\\
Pour quatre raisons. Il n'existe pas de données d'entraînement : prédire le bon type de
solution supposerait des centaines de besoins qualifiés puis validés par leur résultat
réel, corpus dont aucune organisation ne dispose au démarrage. La décision doit être
contestable règle par règle en comité. Le contexte GxP pénalise l'opacité. Et le
déterminisme rend le moteur testable sur des sorties exactes.

\textbf{4. Quelle est la différence entre votre moteur et de l'apprentissage
automatique ?}\\
L'origine de la connaissance. Ici, elle est écrite par un humain et lisible dans le
code ; en apprentissage, elle est inférée statistiquement de données. Il en découle tout
le reste : pas de jeu d'entraînement, déterminisme, explicabilité native — mais aussi
l'incapacité à découvrir ce qu'un expert n'a pas prévu.

\textbf{5. Que ferait un modèle de langage s'il était ajouté ?}\\
Trois choses, toutes en entrée du moteur, jamais à sa place : pré-remplir les quinze
faits de qualification à partir d'une description libre, rédiger un premier jet de fiche
projet, et détecter des doublons sémantiques là où la comparaison lexicale échoue. Le
modèle proposerait des entrées ; les règles continueraient de décider, pour préserver la
traçabilité.

\subsection*{Sur les algorithmes}

\textbf{6. Comment calculez-vous le readiness ?}\\
Chaque dimension porte une note sur 5, un poids et une confiance entre 0 et 1. Le score
est cent fois la somme des notes pondérées par poids et confiance, divisée par la somme
des maxima pondérés de la même façon. La confiance apparaît aux deux termes, donc une
réponse mal étayée s'efface au lieu de fausser le résultat.

\textbf{7. Pourquoi les moyennes d'axe n'utilisent-elles pas la confiance ?}\\
Parce qu'elles décrivent deux choses différentes. Le score global mesure une maturité
consolidée, pour laquelle la solidité du constat compte. La moyenne d'axe décrit le
niveau observé, indépendamment de la qualité de l'instruction.

\textbf{8. Que se passe-t-il si une dimension n'a pas été renseignée ?}\\
Une confiance nulle la marque comme non instruite. Elle sort du calcul du score et des
moyennes — une dimension non mesurée ne vaut pas zéro — et pèse uniquement sur le taux
de complétion. Sans ce mécanisme, ajouter une question au référentiel ferait chuter le
readiness de toutes les campagnes existantes sans qu'aucune réalité n'ait changé.

\textbf{9. Comment le score de priorité est-il construit ?}\\
Six critères notés sur 100, pondérés, moins une pénalité de risque. Les pondérations par
défaut sont 0,25 pour l'impact métier, 0,20 pour l'impact financier, 0,20 pour la
faisabilité, 0,15 pour l'alignement, 0,10 pour le readiness et 0,10 pour les synergies.
La pénalité vaut le produit probabilité par impact, ramené sur 25, multiplié par un
plafond de 15~points.

\textbf{10. Que se passe-t-il si l'utilisateur saisit des poids dont la somme n'est pas
1 ?}\\
Ils sont renormalisés avant application. C'est un choix d'ergonomie : l'utilisateur
pousse le critère qui l'intéresse sans avoir à rééquilibrer les cinq autres à la main,
et les scores restent comparables.

\textbf{11. Pourquoi 15 points de pénalité maximale, et pas 20 ou 10 ?}\\
C'est un étalonnage, pas une constante théorique. Assez élevé pour qu'une opportunité
très risquée soit réellement déclassée, assez bas pour qu'un risque maîtrisable
n'annule pas un impact métier fort. La valeur est isolée dans une constante nommée,
modifiable en un point.

\textbf{12. Pourquoi le payback n'utilise-t-il pas le TCO ?}\\
Parce qu'il répond à une question précise : en combien de mois l'investissement initial
est-il récupéré ? Y intégrer trois années de coûts récurrents non encore engagés
allongerait artificiellement le délai. Le numérateur est donc le CAPEX seul, le
dénominateur le gain annuel net des coûts d'exploitation, ramené au mois.

\textbf{13. Qu'est-ce qu'une règle bloquante dans le moteur de qualification ?}\\
Une condition sans laquelle le type de solution n'a pas de sens — par exemple des
données d'image pour de la vision par ordinateur. Si elle n'est pas satisfaite, le type
est disqualifié et son score divisé par quatre. Il reste visible dans les alternatives
avec le motif de sa disqualification : savoir qu'un type serait pertinent mais qu'il est
bloqué sur la vérité terrain est plus utile qu'une simple absence.

\subsection*{Sur la sécurité}

\textbf{14. Comment garantissez-vous qu'une organisation ne voit pas les données d'une
autre ?}\\
L'identifiant d'organisation est lu dans le jeton signé, jamais dans le corps ou les
paramètres de la requête. Un utilisateur ne peut donc pas se déclarer membre d'une autre
organisation sans forger le jeton, ce qui exigerait la clé de signature. Treize tests
construisent une seconde organisation complète et vérifient qu'aucune lecture, écriture,
suppression ou affectation ne traverse la frontière.

\textbf{15. Pourquoi 404 et non 403 pour une ressource hors périmètre ?}\\
Répondre 403 confirmerait que la ressource existe, ce qui est déjà une fuite. En
revanche, une écriture refusée \emph{à l'intérieur} de sa propre organisation répond
bien 403 : l'utilisateur voit l'objet, il n'a simplement pas le droit d'agir dessus.

\textbf{16. Si je masque un bouton dans React, suis-je protégé ?}\\
Non, et c'est un point que le projet prend au sérieux. Le masquage est une commodité
d'usage. La protection réelle est la dépendance d'autorisation côté serveur. Les tests
appellent les endpoints directement, sans passer par l'interface, précisément parce que
c'est ainsi qu'un attaquant procéderait.

\textbf{17. Un administrateur d'entreprise peut-il devenir administrateur de la
plateforme ?}\\
Non. La permission \code{platform.admin} est la seule à franchir la frontière entre
organisations, et une vérification explicite empêche un administrateur d'entreprise de
se l'attribuer ou de la placer dans un rôle qu'il crée. Deux tests couvrent ces deux
chemins.

\textbf{18. Comment les mots de passe sont-ils stockés ?}\\
Hachés avec bcrypt via passlib. Aucun mot de passe en clair n'est stocké ni journalisé.
Les comptes de démonstration partagent un mot de passe lu d'une variable
d'environnement, jamais écrit dans le code de l'interface.

\textbf{19. Qu'est-ce qui empêche de modifier le journal d'audit ?}\\
L'API n'expose aucune route de mise à jour ni de suppression sur cette table. La
traçabilité repose donc sur l'absence de chemin d'écriture, pas sur une permission qu'on
pourrait accorder par erreur. La table est aussi dépourvue de clés étrangères, pour
qu'une trace survive à la suppression de l'entité qu'elle décrit.

\textbf{20. Quelles faiblesses de sécurité connaissez-vous dans la version actuelle ?}\\
Trois, assumées. La clé de signature et le mot de passe du superutilisateur ont des
valeurs par défaut à redéfinir avant production. Il n'existe ni jeton de
rafraîchissement ni liste de révocation, donc un jeton volé reste exploitable jusqu'à
son expiration. Et il n'y a pas d'authentification à deux facteurs, ce qui serait
attendu en périmètre réglementé.

\subsection*{Sur l'architecture et le code}

\textbf{21. Pourquoi FastAPI plutôt qu'un autre framework ?}\\
Pour trois propriétés utiles ici : la validation des entrées adossée à Pydantic, la
documentation OpenAPI générée automatiquement — précieuse pour démontrer le
comportement de l'API — et le système d'injection de dépendances, qui a permis de
centraliser l'autorisation en deux fonctions plutôt que de la disperser dans 99~routes.

\textbf{22. Pourquoi séparer les services des routers ?}\\
Pour que les formules soient testables sans serveur ni base de données. C'est ce qui a
permis de corriger le calcul du payback et la sémantique du readiness tard dans le
projet sans redouter de régression.

\textbf{23. Comment gérez-vous l'évolution du schéma ?}\\
Par trois migrations Alembic versionnées, écrites pour s'appliquer sur une base déjà
peuplée. La troisième, qui a introduit les périmètres, crée les colonnes avec une valeur
par défaut le temps de la reprise, recopie les rattachements existants, puis retire la
valeur par défaut.

\textbf{24. Pourquoi la table utilisateur s'appelle-t-elle \code{app\_user} ?}\\
Parce que \code{user} est un mot réservé en PostgreSQL.

\textbf{25. Pourquoi conserver \code{department\_id} sur l'utilisateur alors qu'une
table d'association existe ?}\\
La colonne désigne le département principal et sert de repli ; la table d'association
porte les rattachements multiples introduits plus tard. Cette coexistence est le
résultat d'une migration qui devait préserver les données existantes, et non un choix de
conception initial.

\subsection*{Sur le périmètre livré}

\textbf{26. Tous les modules du cahier des charges sont-ils terminés ?}\\
Neuf sur douze le sont. Trois ne le sont que partiellement, et le rapport les détaille :
l'export de la fiche projet produit du HTML imprimable et non un fichier \code{.docx},
le suivi de projet dispose de son API mais pas de son écran dédié, et les notifications
sont persistées sans panneau de consultation.

\textbf{27. Le bouton dit « Export Word / PDF » — est-ce exact ?}\\
Non, et c'est signalé comme tel dans le rapport. L'API renvoie du JSON et une page HTML
mise en page pour l'impression. L'utilisateur obtient un PDF via son navigateur, mais
aucun fichier Word n'est généré. Le libellé est en avance sur la fonction.

\textbf{28. Qu'est-ce qui a été le plus difficile ?}\\
La reprise du contrôle d'accès. La première version ne comptait que dix permissions,
toutes en écriture : les quarante-six endpoints de lecture n'étaient protégés que par
l'authentification, ce qui laissait tout utilisateur connecté lire l'intégralité de son
organisation. Le durcissement a porté le catalogue à 41~permissions et introduit la
notion de périmètre, sans casser les modules existants.

\textbf{29. Comment avez-vous validé que le durcissement ne cassait rien ?}\\
Par une suite de 112~tests, dont 40 couvrent le comportement de chaque rôle et 4 sont
des garde-fous structurels qui parcourent la table de routage pour vérifier qu'aucun
endpoint n'est exposé sans contrôle et qu'aucune dépendance d'autorisation n'a été
silencieusement transformée en paramètre de requête.

\textbf{30. Que feriez-vous différemment en recommençant ?}\\
Je concevrais le modèle de permissions avant les fonctionnalités, et non après. Le
rattraper sur un code déjà écrit a coûté plus cher que de le poser d'emblée, et c'est
la difficulté la plus instructive du projet.

% =====================================================================
%  Annexe B — Références et webographie
% =====================================================================
\chapter{Références et webographie}
\label{ann:references}

Les références ci-dessous correspondent aux documentations effectivement consultées
pendant la réalisation, ainsi qu'aux cadres méthodologiques mentionnés dans le rapport.
Les dates de consultation sont à compléter si l'établissement l'exige.

\section*{Documentations techniques}

\begin{enumerate}[label={[\arabic*]}, leftmargin=*, itemsep=5pt]
  \item \textbf{FastAPI} — documentation officielle du framework, notamment les
    chapitres sur l'injection de dépendances et la sécurité.\\
    \url{https://fastapi.tiangolo.com/}

  \item \textbf{SQLAlchemy 2.0} — guide de l'ORM en style déclaratif typé.\\
    \url{https://docs.sqlalchemy.org/en/20/}

  \item \textbf{Alembic} — documentation des migrations de schéma.\\
    \url{https://alembic.sqlalchemy.org/}

  \item \textbf{Pydantic v2} — validation et sérialisation des données.\\
    \url{https://docs.pydantic.dev/latest/}

  \item \textbf{PostgreSQL} — manuel de référence.\\
    \url{https://www.postgresql.org/docs/}

  \item \textbf{React} — documentation officielle.\\
    \url{https://react.dev/}

  \item \textbf{TypeScript} — manuel du langage.\\
    \url{https://www.typescriptlang.org/docs/}

  \item \textbf{Vite} — documentation de l'outil de build.\\
    \url{https://vite.dev/}

  \item \textbf{React Router} — routage côté client.\\
    \url{https://reactrouter.com/}

  \item \textbf{TanStack Query} — gestion du cache de données serveur.\\
    \url{https://tanstack.com/query/latest}

  \item \textbf{Recharts} — bibliothèque de graphiques pour React.\\
    \url{https://recharts.org/}

  \item \textbf{pytest} — documentation du cadre de test.\\
    \url{https://docs.pytest.org/}

  \item \textbf{Docker Compose} — référence du fichier de composition.\\
    \url{https://docs.docker.com/compose/}
\end{enumerate}

\section*{Normes et spécifications}

\begin{enumerate}[label={[\arabic*]}, leftmargin=*, itemsep=5pt, start=14]
  \item \textbf{RFC 7519} — \emph{JSON Web Token (JWT)}, Internet Engineering Task
    Force, mai 2015. Spécification du format de jeton utilisé pour
    l'authentification.\\
    \url{https://datatracker.ietf.org/doc/html/rfc7519}

  \item \textbf{RFC 6749} — \emph{The OAuth 2.0 Authorization Framework}, IETF, octobre
    2012. Cadre dont FastAPI reprend le flux par mot de passe pour le formulaire de
    connexion.\\
    \url{https://datatracker.ietf.org/doc/html/rfc6749}

  \item \textbf{Spécification OpenAPI} — format de description d'API généré
    automatiquement par l'application sur \code{/docs}.\\
    \url{https://spec.openapis.org/}

  \item \textbf{bcrypt} — Provos, N. et Mazières, D., \emph{A Future-Adaptable Password
    Scheme}, USENIX Annual Technical Conference, 1999. Algorithme de hachage employé
    pour les mots de passe.
\end{enumerate}

\section*{Cadres méthodologiques et référentiels métier}

\begin{enumerate}[label={[\arabic*]}, leftmargin=*, itemsep=5pt, start=18]
  \item \textbf{TOGAF} — \emph{The Open Group Architecture Framework}, The Open Group.
    Cadre dont s'inspirent le découpage du diagnostic AS-IS et le séquencement des
    phases de projet employés dans la plateforme.\\
    \url{https://www.opengroup.org/togaf}

  \item \textbf{Bonnes pratiques GxP} — ensemble des bonnes pratiques applicables à
    l'industrie pharmaceutique (fabrication, laboratoire, distribution), qui fondent
    les garde-fous ajoutés automatiquement par le moteur de qualification lorsqu'une
    opportunité est déclarée en périmètre réglementé.

  \item \textbf{Principes ALCOA+} — exigences d'intégrité des données réglementées
    (attribuable, lisible, contemporain, original, exact, complet, cohérent, durable,
    disponible). Ces principes justifient le caractère immuable du journal d'audit.
\end{enumerate}

\begin{remarque}[Sur la portée de cette liste]
Cette webographie recense les sources techniques réellement utilisées. Les références
méthodologiques TOGAF, GxP et ALCOA+ ont servi de cadre conceptuel au découpage
fonctionnel ; la plateforme s'en inspire sans prétendre à une conformité certifiée,
laquelle relèverait d'une démarche de validation formelle non menée dans ce projet.
\end{remarque}


\end{document}
